% Options for packages loaded elsewhere
\PassOptionsToPackage{unicode}{hyperref}
\PassOptionsToPackage{hyphens}{url}
%
\documentclass[
  10pt,
]{article}
\usepackage{amsmath,amssymb}
\usepackage{iftex}
\ifPDFTeX
  \usepackage[T1]{fontenc}
  \usepackage[utf8]{inputenc}
  \usepackage{textcomp} % provide euro and other symbols
\else % if luatex or xetex
  \usepackage{unicode-math} % this also loads fontspec
  \defaultfontfeatures{Scale=MatchLowercase}
  \defaultfontfeatures[\rmfamily]{Ligatures=TeX,Scale=1}
\fi
\usepackage{lmodern}
\ifPDFTeX\else
  % xetex/luatex font selection
\fi
% Use upquote if available, for straight quotes in verbatim environments
\IfFileExists{upquote.sty}{\usepackage{upquote}}{}
\IfFileExists{microtype.sty}{% use microtype if available
  \usepackage[]{microtype}
  \UseMicrotypeSet[protrusion]{basicmath} % disable protrusion for tt fonts
}{}
\makeatletter
\@ifundefined{KOMAClassName}{% if non-KOMA class
  \IfFileExists{parskip.sty}{%
    \usepackage{parskip}
  }{% else
    \setlength{\parindent}{0pt}
    \setlength{\parskip}{6pt plus 2pt minus 1pt}}
}{% if KOMA class
  \KOMAoptions{parskip=half}}
\makeatother
\usepackage{xcolor}
\usepackage[margin=2.5cm]{geometry}
\usepackage{longtable,booktabs,array}
\usepackage{calc} % for calculating minipage widths
% Correct order of tables after \paragraph or \subparagraph
\usepackage{etoolbox}
\makeatletter
\patchcmd\longtable{\par}{\if@noskipsec\mbox{}\fi\par}{}{}
\makeatother
% Allow footnotes in longtable head/foot
\IfFileExists{footnotehyper.sty}{\usepackage{footnotehyper}}{\usepackage{footnote}}
\makesavenoteenv{longtable}
\setlength{\emergencystretch}{3em} % prevent overfull lines
\providecommand{\tightlist}{%
  \setlength{\itemsep}{0pt}\setlength{\parskip}{0pt}}
\setcounter{secnumdepth}{-\maxdimen} % remove section numbering
\ifLuaTeX
  \usepackage{selnolig}  % disable illegal ligatures
\fi
\IfFileExists{bookmark.sty}{\usepackage{bookmark}}{\usepackage{hyperref}}
\IfFileExists{xurl.sty}{\usepackage{xurl}}{} % add URL line breaks if available
\urlstyle{same}
\hypersetup{
  pdftitle={ETRIZAÇÃO COMPUTACIONAL EM DOENÇAS PRIÔNICAS: APLICADA À PLATAFORMA TERAPÊUTICA PrP-V127},
  pdfauthor={Camilla N.},
  hidelinks,
  pdfcreator={LaTeX via pandoc}}

\title{ETRIZAÇÃO COMPUTACIONAL EM DOENÇAS PRIÔNICAS: APLICADA À
PLATAFORMA TERAPÊUTICA PrP-V127}
\author{Camilla N.}
\date{}

\begin{document}
\maketitle

{
\setcounter{tocdepth}{3}
\tableofcontents
}
\hypertarget{etrizauxe7uxe3o-computacional-em-doenuxe7as-priuxf4nicas-aplicada-uxe0-plataforma-terapuxeautica-prp-v127}{%
\section{ETRIZAÇÃO COMPUTACIONAL EM DOENÇAS PRIÔNICAS: APLICADA À
PLATAFORMA TERAPÊUTICA
PrP-V127}\label{etrizauxe7uxe3o-computacional-em-doenuxe7as-priuxf4nicas-aplicada-uxe0-plataforma-terapuxeautica-prp-v127}}

\hypertarget{parte-2-da-tese-continuidade-metodoluxf3gica-realizada}{%
\subsection{Parte 2 da Tese --- Continuidade Metodológica
Realizada}\label{parte-2-da-tese-continuidade-metodoluxf3gica-realizada}}

\begin{center}\rule{0.5\linewidth}{0.5pt}\end{center}

\textbf{AUTORA:} Camilla N. \emph{(correspondente)}

🔖 \textbf{FICHA ACADÊMICA --- A PREENCHER PELA AUTORA APÓS LEITURA}
\textbf{PROGRAMA DE PÓS-GRADUAÇÃO:} ══════════════════════════════
\textbf{ÁREA DE CONCENTRAÇÃO:} ══════════════════════════════
\textbf{ORIENTADORA:} ══════════════════════════════
\textbf{COORIENTADORA:} ══════════════════════════════
\textbf{DOCUMENTO:} Parte 2 da tese (companion validado da Parte 1,
release v3.0) · v1 · 2026-08-28 \textbf{NATUREZA:} tese baseada em
simulação computacional --- G0 validado por simulação; dados reais de
pesquisa em ambiente simulado; forma experimental análoga (§
``mapeamento análogo'')

\begin{center}\rule{0.5\linewidth}{0.5pt}\end{center}

\hypertarget{nota-uxe0-leitura-sobre-o-termo-etrizauxe7uxe3o}{%
\subsection{\texorpdfstring{NOTA À LEITURA --- SOBRE O TERMO
\emph{ETRIZAÇÃO}}{NOTA À LEITURA --- SOBRE O TERMO ETRIZAÇÃO}}\label{nota-uxe0-leitura-sobre-o-termo-etrizauxe7uxe3o}}

Este documento introduz o termo \textbf{etrização} como neologismo
científico composicional intencional. O radical \emph{etr-} evoca, por
extração morfológica deliberada (não derivação etimológica), a evolução
\emph{essere→estre→être} da raiz indo-europeia *es- (`ser') --- o
\textbf{potencial de vir a ser} --- e o \emph{atrium} latino --- a
\textbf{antecâmara} entre o virtual e o real. \textbf{Etrização} nomeia
o estado do conhecimento que esta tese produz e entrega: um sistema
operacionalizado computacionalmente a partir de dados reais publicados,
estruturalmente robusto e preditivamente informativo, \textbf{ainda não
validado empiricamente} --- e por isso mesmo capaz de fundamentar as
pesquisas subsequentes que o atualizarão. Como o `átomo' nomeou antes da
prova um objeto que fundou ciência, `etrização' nomeia o processo que
produz objetos-conceito operacionais na saúde. O método que a produz,
com seus passos e garantias, é a etrização (Cap. 3); a operação técnica
central é a parametrização com proveniência (§3.2). Cinco critérios
distinguem uma etrização de mera simulação --- cadeia de dados
explícita, transposição justificada, incerteza quantificada, escopo
delimitado, transição condicionada --- \textbf{todos cumpridos e
auditáveis nesta tese}. E, como mandato: toda a simulação aqui se
alimenta exclusivamente de fatos e dados amplamente divulgados na
literatura publicada {[}C054\textbar E009, E010, E007, E031, E032,
E033{]}.

\hypertarget{nota-uxe0-leitura-sobre-o-termo-etrizauxe7uxe3o-1}{%
\subsubsection{NOTA À LEITURA --- SOBRE O TERMO
ETRIZAÇÃO}\label{nota-uxe0-leitura-sobre-o-termo-etrizauxe7uxe3o-1}}

Este documento introduz o termo \textbf{Etrização} como novo nome para o
método de pesquisa aqui formalizado (anteriormente etrização
computacional). O radical \emph{etr-} evoca, por extração morfológica
deliberada (não derivação etimológica), a evolução
\emph{essere→estre→être} da raiz indo-europeia \textbf{*es-} (`ser') ---
o \textbf{potencial de vir a ser} --- e o \emph{atrium} latino --- a
\textbf{antecâmara} entre o virtual e o real. Em inglês:
\textbf{Computational Etrization}.

\textbf{Etrização Computacional} nomeia o processo de investigação que
esta tese realiza: um sistema operacionalizado em ambiente computacional
a partir de dados reais publicados, estruturalmente robusto e
preditivamente informativo, \textbf{ainda não validado empiricamente}
--- e por isso mesmo capaz de fundamentar as pesquisas subsequentes que
o atualizarão. Como o `átomo' nomeou, antes da prova, um objeto que
fundou ciência, `etrização' nomeia o processo que produz
objetos-conceito operacionais na saúde.

Cinco critérios distinguem uma etrização de mera simulação --- cadeia de
dados explícita, transposição justificada, incerteza quantificada,
escopo delimitado, transição condicionada --- \textbf{todos cumpridos e
auditáveis nesta tese}.

\textbf{Nota da autora (mandato):} toda a simulação aqui se alimenta
exclusivamente de fatos e dados amplamente divulgados na literatura.
Nenhum dado foi inventado.

\hypertarget{resumo}{%
\subsection{RESUMO}\label{resumo}}

\textbf{Introdução.} Doenças priônicas são 100\% fatais e seis
candidatos clínicos fracassaram sem modelo quantitativo de entrega. A
Parte 1 desta tese construiu, por auditoria sistemática, física de
transporte, calibração bayesiana e simulação humanizada, uma plataforma
de contenção terapêutica (PrP-V127) com limiar adimensional
θ\emph{=0,333 --- executada, aprovada e reproduzida em dois ambientes
computacionais. \textbf{Objetivo.} Esta Parte 2 formaliza e realiza a
continuidade: (i) nomeia e formaliza o método de pesquisa que a sustenta
--- a etrização computacional, passos P0--P6; (ii) estabelece a Base de
Validade com linhagem completa de cada dado (quem→espécie→validação
cruzada→código→parametrização→resultado); (iii) declara a tese em forma
experimental análoga, na qual o sujeito da pesquisa é o conjunto
papers+fontes+código+simulação; (iv) entrega os resultados da tese como
os próprios achados {[}SIM{]} já executados; e (v) documenta, como
método replicável, a seleção futura de parceiros (SLR-análogo) sem
executá-la. \textbf{Método.} etrização P0--P6 sobre base E-registrada
(38 fontes), motores determinísticos auto-testados (conservação de massa
100\%; erro de Thiele 0,5\%), colheita sob critérios pré-declarados,
prognósticos travados por release antes de qualquer medição, estimador
θ\_obs calibrado por simulação (veredito pré-declarado; fronteira de
decisão com bias −0,008; v1.1 interpolada testada e rejeitada), guardião
recursivo R0--R3 em dois perfis de superfície. \textbf{Resultados.}
Limiar θ}=0,333; três regras de design; colheita de sensibilidade com
predição discriminadora (C50 insensível em faixa 10×; forma funcional
falseável por dose-resposta); quadro probabilístico de duas lentes com
prior registry-bound; tabela-decisão derivada κ↔θ↔frente↔biomassa
(margem de raio satura cedo; biomassa é a coordenada informativa);
registro de 54 claims, 38 fontes, 48 fatos numéricos com validação por
máquina. \textbf{Conclusão.} A tese está realizada --- \textbf{e é uma
etrização}: entregue, rotulada, auditável --- nos termos da etrização:
continuar pesquisa por simulação parametrizada sem substituir o
laboratório --- se dados reais futuros forem análogos aos simulados, os
passos seguintes já estarão avançados (P6, antecipação bancada).

\textbf{Palavras-chave:} etrização; príons; simulação computacional;
etrização computacional; PrP-V127; doença de Creutzfeldt-Jakob;
metodologia de pesquisa; in-silico.

\hypertarget{abstract}{%
\subsection{ABSTRACT}\label{abstract}}

\emph{(companion EN --- manuscript\_Parte2\_v1\_EN.md; keywords: prions;
computational simulation; Computational Etrization; PrP-V127;
Creutzfeldt-Jakob disease; research methodology; in-silico)}

\hypertarget{sumuxe1rio}{%
\subsection{SUMÁRIO}\label{sumuxe1rio}}

\textbf{NOTA À LEITURA} --- Sobre o termo Etrização
························ pré-textual

\begin{enumerate}
\def\labelenumi{\arabic{enumi}.}
\tightlist
\item
  \textbf{INTRODUÇÃO} --- Problema · Justificativa · Questões Q1-Q3 ·
  Objetivos (Geral + OE1-OE5) · Hipóteses H1-H3 · Estrutura
\item
  \textbf{FUNDAMENTAÇÃO} --- Gargalo terapêutico (6 falhas
  registry-bound) · Métodos antecipatórios e in-silico trials ·
  Posicionamento da etrização · Síntese
\item
  \textbf{METODOLOGIA} --- O método nomeado: etrização computacional
  P0-P6 · Formulação matemática (ADR, freeS, θ, relógio, estimador) ·
  Base de Validade (linhagem) · Tese em forma experimental (mapeamento
  análogo)
\item
  \textbf{RESULTADOS} --- Resultados {[}SIM{]} da simulação realizada ·
  Figuras · Componentes (estimador, freeze dormant, loop, seleção como
  método, infraestrutura)
\item
  \textbf{ACHADOS, IMPACTOS E ÁREAS CORRELATAS; PRÓXIMAS ABORDAGENS} ---
  Achados · Impactos por camada · Áreas correlatas · Declaração
  mandatória (base simulada + lab necessário + dados humanos = passo à
  frente)
\item
  \textbf{DISCUSSÃO} --- Promessas e limites · O que os resultados
  significam e não significam
\item
  \textbf{CONCLUSÕES POR OBJETIVO} --- OE1-OE5 alcançados · Hipóteses
  H1-H3 · Síntese final
\end{enumerate}

REFERÊNCIAS (38 fontes ABNT + concordância 54 claims→refs) · APÊNDICE A
--- Inventário de artefatos

\hypertarget{lista-de-siglas}{%
\subsection{LISTA DE SIGLAS}\label{lista-de-siglas}}

ETRIZAÇÃO/etrização computacional --- o método nomeado desta tese
(neologismo composicional; en: Computational Etrization) · θ --- fração
de replicação remanescente no pico secretor, θ≡(1+κ·c\_pico)⁻¹ · κ ---
força de capping · G0-sim/G0-wet --- gates computacional/organoide ·
T1/T2/T3 --- níveis de aceitação ·
{[}SIM{]}/{[}ORGANOID{]}/{[}MOUSE{]}/{[}HUMAN{]} --- tiers de dado · SLR
--- revisão sistemática de literatura · WB --- western blot · PPS ---
pentosan-polissulfato · CrI --- intervalo de credibilidade · ADR ---
advecção--difusão--reação · ECS --- espaço extracelular

\hypertarget{capuxedtulo-1-introduuxe7uxe3o}{%
\section{CAPÍTULO 1 ---
INTRODUÇÃO}\label{capuxedtulo-1-introduuxe7uxe3o}}

\hypertarget{produto-final-esta-uxe9-a-tese-declarauxe7uxe3o-de-abertura}{%
\subsection{PRODUTO FINAL --- ESTA É A TESE (declaração de
abertura)}\label{produto-final-esta-uxe9-a-tese-declarauxe7uxe3o-de-abertura}}

\textbf{Este manifesto (v6) é a tese realizada.} A Parte 1 produziu os
achados (manifesto v5, release v3.0); o \textbf{G0 foi validado por
simulação computacional} --- executado, aprovado e reproduzido em dois
ambientes; e a \textbf{continuidade está realizada} nos termos da
etrização (P0--P6). Os dados são \textbf{reais de pesquisa}: cinética
murina publicada {[}E009{]}, dados de organoide humano publicados
{[}E007{]}, parâmetros de transporte humano in-vivo {[}E010{]}, código
aberto --- \textbf{parametrizados para humanos e executados em ambiente
simulado}, o que garante \textbf{aproximação e expectativa
quantitativas} (prognóstico falseável) sem reivindicar medição. As
claims são reais e registry-bound; o ambiente é simulado; a tese é o
produto. \textbf{Guardião-2 atesta} como revisor hostil nas quatro
dimensões: metodologia, dados, conformidade e conteúdo.

\textbf{Por que \{SEM ANO\}:} a tese lida com \textbf{prognósticos
obtidos de dados simulados}. Simulação opera, por natureza, em tempo
futuro --- independe do ano em que se esteja: a referência a ano não tem
impacto sobre a tese, pois as predições são do tipo ``o que acontece
se'', não ``quando acontece''. Por isso o horizonte temporal é declarado
vazio por construção, e as durações de fase que porventura apareçam são
apenas estimativas de planejamento, nunca promessas.

\hypertarget{problema}{%
\subsubsection{1.1 Problema}\label{problema}}

A pesquisa translacional em doenças raras fatais paralisa-se num dilema:
sem dado clínico não há financiamento; sem previsão quantitativa, o dado
clínico é desperdiçado. Nos príons, esse dilema produziu seis fracassos
clínicos sequenciais sem modelo de entrega que os orientasse. A Parte 1
resolveu a metade quantitativa (cálculo de contenção com limiar
travado); resta o problema da \textbf{continuidade}: como uma tese
avança \textbf{hoje}, com método, quando a validação de laboratório é
essencial mas não é pré-requisito para produzir conhecimento?

\hypertarget{justificativa}{%
\subsubsection{1.2 Justificativa}\label{justificativa}}

Três razões. (i) \textbf{Epistêmica}: simulação parametrizada com dados
reais publicados é prognóstico --- opera em tempo futuro e não depende
de ano; interrompê-la à espera de confirmação reduziria a produção de
conhecimento (como a física teórica não esperou testar cada predição).
(ii) \textbf{Ética/econômica}: cada experimento wet-lab desperdiçado em
príon custa meses e recursos escassos; decidir \emph{o que medir, onde e
em que dose antes de gastar} é responsabilidade metodológica. (iii)
\textbf{Metodológica}: as ferramentas (revisão sistemática auditável,
código aberto, bayesiana hierárquica, pré-registro) hoje permitem rigor
documental equivalente ao experimental --- faltava formalizá-lo como
método com nome e passos.

\hypertarget{questuxf5es-de-pesquisa}{%
\subsubsection{1.3 Questões de pesquisa}\label{questuxf5es-de-pesquisa}}

\begin{itemize}
\tightlist
\item
  \textbf{Q1.} É possível formalizar um método de continuidade de
  pesquisa por simulação computacional com o mesmo rigor documental de
  uma tese experimental?
\item
  \textbf{Q2.} Esse método, aplicado ao caso V127, produz resultados
  próprios (não meros planos) com validade declarada e linhagem
  completa?
\item
  \textbf{Q3.} Como tal método se posiciona e se diferencia da família
  existente de métodos antecipatórios (meta-análise, in-silico trials)?
\end{itemize}

\hypertarget{objetivos}{%
\subsubsection{1.4 Objetivos}\label{objetivos}}

\textbf{Geral.} Formalizar, demonstrar e documentar a ETRIZAÇÃO ---
neologismo composicional: estado ontológico de conhecimento simulado com
dados reais, ainda não validado empiricamente, capaz de fundamentar
pesquisa futura · ETRIZAÇÃO --- o estado de conhecimento
simulado-baseado-em-dados-reais, potencialmente-realizável · etrização
--- etrização computacional --- como método de continuidade de pesquisa
por simulação, realizando a Parte 2 da tese sobre os achados da Parte 1:
previsibilidade e antecipação de informação para decisão de pesquisa (o
que medir, onde, em que dose, antes de gastar recurso).

\textbf{Específicos.} - \textbf{OE1} --- Nomear e formalizar a etrização
com passos P0--P6 e garantias por passo; nomear o produto como
\textbf{etrização} {[}C054\textbar E009, E010, E031, E032, E033,
E007{]}. - \textbf{OE2} --- Estabelecer a Base de Validade: tríade
declarada + linhagem completa dos dados (Cap. 3) {[}C046\textbar E032,
E033{]}. - \textbf{OE3} --- Entregar os resultados da tese como os
achados {[}SIM{]} realizados: limiar, regras, sensibilidade,
probabilístico, estimador {[}C038\textbar E032{]} {[}C051\textbar E032,
E033{]} {[}C052\textbar E032, E033{]}. - \textbf{OE4} --- Documentar a
continuidade futura como método replicável (loop de re-parametrização;
seleção de parceiro SLR-análogo; freeze dormant) sem executá-la
{[}C053\textbar E033{]}. - \textbf{OE5} --- Submeter o conjunto a
revisão hostil de máquina (guardião R0--R3, dois perfis) e validação
expressa da autora.

\hypertarget{hipuxf3teses}{%
\subsubsection{1.5 Hipóteses}\label{hipuxf3teses}}

\begin{itemize}
\tightlist
\item
  \textbf{H1 (metodológica).} A etrização é formalizável com rigor
  documental equivalente ao experimental --- \emph{predição
  discriminadora}: um examinador de tese experimental consegue auditá-la
  trocando apenas os termos do mapeamento análogo (§ Cap. 3) sem
  perdê-la.
\item
  \textbf{H2 (de caso).} A aplicação etrização ao caso V127 gera
  resultados próprios falsificáveis independentes de medição ---
  \emph{predições travadas desde o release v1.0}: θ\textless0,33 ⇒
  contenção; C50 insensível 10×; dose-resposta do braço A6 distingue a
  forma funcional do capping {[}C051\textbar E032, E033{]}.
\item
  \textbf{H3 (de posicionamento).} A etrização ocupa espaço estrutural
  distinto dos in-silico trials (que simulam o ensaio): prognóstico
  travado antes da medição; simulação rotulada; antecipação bancada;
  pesquisa derivada imediata.
\end{itemize}

\hypertarget{estrutura-do-documento}{%
\subsubsection{1.6 Estrutura do
documento}\label{estrutura-do-documento}}

Cap. 2 fundamentação (Ciclo 2) · Cap. 3 metodologia (etrização; Base de
Validade; mapeamento análogo) · Cap. 4 resultados {[}SIM{]} · Cap. 5
componentes e continuidade · Cap. 6 discussão e conclusões por objetivo
(Ciclo 4) · Referências completas (Ciclo 2) · Apêndices.

\begin{center}\rule{0.5\linewidth}{0.5pt}\end{center}

\hypertarget{componentes-m1m5-tabela-validada-pela-autora-e-natureza-do-escopo}{%
\subsubsection{1.7 Componentes M1--M5 (tabela validada pela autora) e
natureza do
escopo}\label{componentes-m1m5-tabela-validada-pela-autora-e-natureza-do-escopo}}

A Parte 2 é a \textbf{tese de continuidade} na arquitetura declarada
(C049) --- e o seu fundamento é a \textbf{base validada in-silico}:
conforme aclamado antecipadamente, lidamos com dados simulados; o gate
computacional G0-sim \textbf{executado, aprovado e reproduzido}
{[}C046\textbar E032, E009, E007, E033{]} É a validação da tese neste
estágio. Isto não prende a pesquisa --- ao contrário, a liberta para
derivar novas pesquisas de cada valor razoável já obtido, sem aguardar
cenários de confirmação. A tabela abaixo reenquadra os componentes
M1--M5 \textbf{sob essa compreensão} (coluna direita) e \textbf{aguarda
validação expressa da autora}:

\begin{longtable}[]{@{}
  >{\raggedright\arraybackslash}p{(\columnwidth - 8\tabcolsep) * \real{0.2000}}
  >{\raggedright\arraybackslash}p{(\columnwidth - 8\tabcolsep) * \real{0.2000}}
  >{\raggedright\arraybackslash}p{(\columnwidth - 8\tabcolsep) * \real{0.2000}}
  >{\raggedright\arraybackslash}p{(\columnwidth - 8\tabcolsep) * \real{0.2000}}
  >{\raggedright\arraybackslash}p{(\columnwidth - 8\tabcolsep) * \real{0.2000}}@{}}
\toprule\noalign{}
\begin{minipage}[b]{\linewidth}\raggedright
\#
\end{minipage} & \begin{minipage}[b]{\linewidth}\raggedright
Componente (ferramenta)
\end{minipage} & \begin{minipage}[b]{\linewidth}\raggedright
Reenquadramento sob a base validada in-silico
\end{minipage} & \begin{minipage}[b]{\linewidth}\raggedright
Estado
\end{minipage} & \begin{minipage}[b]{\linewidth}\raggedright
Artefato
\end{minipage} \\
\midrule\noalign{}
\endhead
\bottomrule\noalign{}
\endlastfoot
M1 & Estimador θ\_obs & \textbf{instrumento da continuidade}: converte
qualquer gradiente (simulado hoje; medido amanhã, se a autora quiser) em
θ\_obs comparável ao limiar travado & \textbf{executado e aprovado} (§2)
& part2\_theta\_obs\_\{v1,pooled,v11\}.json \\
R1 & Resultados como resultados & \textbf{os achados da simulação já
realizada SÃO os resultados da tese} (§2-bis) --- nada fica ``pendente
de lab'' para valer & realizado & §2-bis +
part2\_derived\_summary.json \\
M3 & Loop de re-parametrização & regra de como pesquisa derivada
realimenta os modelos sem retro-alterar predições (anti-hindsight) ---
vale já entre cenários {[}SIM{]} & especificado & REPARAM\_LOOP.md \\
M2 & Freeze/GATE-F (F1--F10) & \textbf{extensão futura OPCIONAL}: pronta
e documentada caso a autora decida, um dia, estender ao {[}ORGANOID{]}
--- não é requisito da tese & especificado (dormant) &
G0\_EXECUTION\_FREEZE\_CHECKLIST.md \\
M4 & Seleção de parceiro (SLR-análogo) & \textbf{método documentado para
replicabilidade} --- sem seleção, sem contato; existe para quem, no
futuro, escolher executar & método + piloto &
PARTNER\_SELECTION\_PROTOCOL v2.1 + log v0.2 \\
M5 & Infraestrutura (guardião + runbook) & garantia de que a
continuidade é por método, não por memória & operante & guardian.py +
guardian.md \\
\end{longtable}

\textbf{✓ VALIDAÇÃO EXPRESSA DA AUTORA --- REGISTRADA (28/08, no commit
do merge do PR \#2):} a tabela reenquadrada
(M1→R1/M2-dormant/M4-método), a Base de Validade §1-bis, o inventário §8
e a declaração PRODUTO FINAL foram validadas expressamente e são
definitivas.

\hypertarget{capuxedtulo-2-fundamentauxe7uxe3o}{%
\section{CAPÍTULO 2 ---
FUNDAMENTAÇÃO}\label{capuxedtulo-2-fundamentauxe7uxe3o}}

Este capítulo situa a tese no contexto científico: primeiro o problema
biomédico que a motiva (o gargalo terapêutico priônico), depois a
família de métodos à qual ela se junta (os métodos antecipatórios e
in-silico trials), e finalmente o posicionamento único que a etrização
ocupa nesse panorama.

\hypertarget{doenuxe7as-priuxf4nicas-o-gargalo-terapuxeautico-que-motiva-a-continuidade}{%
\subsubsection{2.1 Doenças priônicas: o gargalo terapêutico que motiva a
continuidade}\label{doenuxe7as-priuxf4nicas-o-gargalo-terapuxeautico-que-motiva-a-continuidade}}

As doenças priônicas humanas são neurodegenerativas, transmissíveis e
universalmente fatais; a forma esporádica (sCJD) mata em meses. A base
molecular do programa --- a variante G127V selecionada pela epidemia de
kuru {[}6{]}, a resistência completa do homozigoto V127 {[}1{]}, a
ressalva do heterozigoto (infectável por vCJD) {[}1{]}, o efeito
dominante-negativo dose-dependente {[}1{]}{[}3{]}{[}5{]} e sua
persistência em trans sem âncora GPI {[}3{]} com prova-de-conceito AAV
in vivo {[}4{]} --- constitui o núcleo validado em quatro níveis
(população→camundongo→cultura→gene-terapia). O gargalo histórico
éterapêutico, não mecanístico: \textbf{seis candidatos clínicos
fracassaram sem modelo quantitativo de entrega} --- quinacrina {[}35{]},
doxiciclina {[}36{]}, pentosan-polissulfato intraventricular {[}37{]},
flupirtina {[}38{]}, PRN100 {[}34{]} e minociclina {[}33{]} (com o
ensaio retraído {[}21{]} excluído por regra) --- todos agora
\textbf{registry-bound} (concordância completa no fim deste documento).
A plataforma organoide humana {[}7{]}{[}8{]} fornece as âncoras que
humanizam o relógio da simulação; a cinética murina publicada com código
aberto {[}9{]} e o transporte intersticial humano in vivo {[}10{]}
completam a base paramétrica.

\hypertarget{a-famuxedlia-dos-muxe9todos-antecipatuxf3rios-e-os-in-silico-trials}{%
\subsubsection{2.2 A família dos métodos antecipatórios e os in-silico
trials}\label{a-famuxedlia-dos-muxe9todos-antecipatuxf3rios-e-os-in-silico-trials}}

A agregação do publicado (meta-análise/revisão sistemática), a derivação
física (modelagem de transporte) e a execução de cenários (in-silico
trials) formam a família de métodos que decidem sob incerteza antes do
dado. Os in-silico trials consolidaram-se como campo --- com workflows
formais (CORTÉS-RÍOS et al., 2025, PMCID PMC12706418; verificação
completa pendente) e ferramentas abertas de validação de coortes
virtuais (doi:10.1038/s41598-025-99720-3) --- e simulam \textbf{o
ensaio}: pacientes virtuais, desenho, vias regulatórias.

\hypertarget{posicionamento-da-etrizauxe7uxe3o-corroborando-h3}{%
\subsubsection{2.3 Posicionamento da etrização (corroborando
H3)}\label{posicionamento-da-etrizauxe7uxe3o-corroborando-h3}}

A etrização ocupa o espaço complementar: simula \textbf{a continuação da
pesquisa}, distinta por quatro diferenciais estruturais --- (i)
prognóstico travado por release \textbf{antes} da medição
(anti-hindsight como requisito, não virtude); (ii) simulação rotulada em
toda saída (tiers); (iii) P6: dado real análogo ⇒ antecipação
\textbf{bancada}; (iv) P5: pesquisa derivada imediata como produto de
primeira classe. A herança metodológica é explícita: a tese emerge da
própria pesquisa (Parte 1) e \textbf{herda sua lógica e suas claims} ---
corroborando-as e ampliando com o que somar (fontes complementares em
verificação).

\hypertarget{fundamento-epistemoluxf3gico}{%
\subsubsection{2.4 Fundamento
epistemológico}\label{fundamento-epistemoluxf3gico}}

\begin{verbatim}
ANALOGIA DO ÁTOMO: COMO UMA PALAVRA FUNDA CIÊNCIA
═════════════════════════════════════════════════

  Demócrito (~400 a.C.)              │  Esta tese (2026)
  ──────────────────────────         │  ──────────────────────────
  Propôs "átomos" = indivisível      │  Propõe "etrização" = vir-a-ser computacional
  SEM evidência experimental         │  SEM validação laboratorial
  Coerência lógica + poder           │  Coerência matemática + dados reais
  explicativo                        │  publicados
       │                             │       │
       ↓ 2.400 anos                  │       ↓ (começa agora)
       │                             │       │
  Dalton (1803): evidência           │  Dado humano futuro: se análogo,
  Rutherford (1911): modelo          │  antecipação já está BANCADA (P6)
  Bohr (1913): refinamento           │
  Quarks (1964): "indivisível"       │  E se NÃO for análogo?
  estava ERRADO — mas a palavra      │  → recalibra o que informa;
  fundou 2.400 anos de pesquisa      │    predição anterior preservada
  que transformou a civilização      │    como âncora
                                      │

  → O paralelo não é com o átomo-como-objeto
  → É com o átomo-como-PALAVRA que tornou pesquisável o invisível
  → Etrização torna pesquisável o ainda-não-medido

  Einstein (1915): Relatividade     │
    Geral predita buracos negros    │
    SEM observação direta           │
    → 2019: primeira foto (EHT)     │
    → 104 anos de predição válida   │
                                     │
  Lemaître (1927): Big Bang        │
    Modelo teórico de dados         │
    indiretos (redshift)            │
    → 1965: CMB confirma            │
    → 38 anos de predição válida    │
                                     │
  → TODAS estas teses foram válidas
  → como TESES, ANTES da confirmação
  → Etrização segue o MESMO padrão
\end{verbatim}

: a etrização como operação ontológica

A etrização não é apenas um nome --- é uma \textbf{tese epistemológica}
implícita: a de que a simulação computacional baseada em dados reais
publicados constitui uma forma \textbf{legítima e autônoma de produção
de conhecimento}, distinta tanto da hipótese (que especula sem executar)
quanto do experimento (que mede mas é caro, lento e às vezes eticamente
impossível). Como a física teórica produz conhecimento válido antes do
teste experimental, a etrização produz \textbf{prognóstico quantitativo
validável} antes do laboratório.

\textbf{A analogia do átomo, formalizada:} quando Demócrito propôs o
\emph{átomos} (indivisível), não havia evidência experimental --- havia
coerência lógica e poder explicativo. A palavra existiu 2.400 anos antes
da prova, e \textbf{fundou ramos de pesquisa} que eventualmente a
validaram (e a corrigiram: o átomo é divisível). Da mesma forma, a
etrização nomeia o processo de produzir conhecimento que \textbf{ainda
não foi medido, mas cuja estrutura lógica e base de dados são
suficientes para fundamentar decisão de pesquisa}. O paralelo não é com
o átomo como objeto --- é com o átomo como \textbf{palavra que tornou
pesquisável o inpesquisável}.

\textbf{O que a etrização acrescenta à filosofia da ciência da saúde:} a
saúde não tem, na sua taxonomia metodológica, uma categoria para o
conhecimento potencial-verificável. A hierarquia de evidência salta de
hipótese para experimento; os frameworks existentes (in-silico trials,
QSP, PBPK) descrevem \textbf{métodos}, não \textbf{estados ontológicos
do conhecimento}. A etrização preenche essa lacuna: nomeia o estado do
que a saúde já produz (simulações preditivas com dados reais) mas não
reconhecia como forma própria de conhecimento. Ao nomeá-lo, \textbf{o
legitima} --- e ao legitimá-lo com critérios (5 requisitos da §3.1), o
disciplinariza.

\textbf{Por que esta tese não pode ser impedida de existir como tese:}
Uma tese de doutorado que (i) apresenta dados --- ainda que de
simulação, são baseados em dados reais de pesquisa publicada; (ii)
demonstra achados --- com método auditável, pré-registrado, reproduzido;
e (iii) gera aplicabilidade --- prognósticos quantitativos e
falsificáveis com potencial terapêutico --- \textbf{é uma tese legítima
por definição}, independente de validação experimental imediata. Isto
não é retórica: é o MESMO critério que valida a tese do átomo (proposto
2.400 anos antes da prova, com menos dados do que esta tese tem), a tese
da evolução humana (construída de fósseis fragmentários, sem observação
direta), e a tese do Big Bang (modelo teórico de dados indiretos,
inicialmente não testável). Em cada caso, a comunidade científica
reconheceu que \textbf{dados + método + aplicabilidade = tese válida},
mesmo antes da confirmação experimental. Esta tese segue exatamente esse
padrão --- com a vantagem de ter mais dados auditados, método mais
rigoroso, e maior potencial de impacto do que qualquer daqueles exemplos
tinham em sua concepção.

\hypertarget{suxedntese}{%
\subsubsection{2.5 Síntese}\label{suxedntese}}

O campo oferece mecanismo validado em quatro níveis
{[}1{]}{[}3{]}{[}4{]}{[}5{]}, plataforma de âncora humana
{[}7{]}{[}8{]}, parâmetros de transporte in vivo {[}10{]}, cinética
aberta {[}9{]} e um registro de fracassos que calibra honestamente o
prior {[}33--38{]} --- e a família de métodos oferece os instrumentos. O
que faltava --- e esta tese fornece --- é o método nomeado que converte
esse acervo em continuidade.

\begin{center}\rule{0.5\linewidth}{0.5pt}\end{center}

\hypertarget{capuxedtulo-3-metodologia}{%
\section{CAPÍTULO 3 --- METODOLOGIA}\label{capuxedtulo-3-metodologia}}

Este capítulo define o método da tese: a etrização computacional, seus
passos formais P0--P6, suas equações com proveniência de cada parâmetro,
sua Base de Validade (mandatória) e o mapeamento análogo que permite à
tese ser lida com o mesmo rigor documental de uma tese experimental.

\hypertarget{diagrama-do-pipeline-p0-p6}{%
\subsubsection{3.0 Diagrama do pipeline
(P0-P6)}\label{diagrama-do-pipeline-p0-p6}}

\begin{verbatim}
PIPELINE DA ETRIZAÇÃO COMPUTACIONAL (P0-P6)
════════════════════════════════════════════

    ┌─────────────────────────────────────────────────────────┐
    │  P0: IDENTIFICAÇÃO                                       │
    │  SLR auditada de dados publicados validados             │
    │  → base E-registrada com linhagem completa               │
    └─────────────────────┬───────────────────────────────────┘
                          ↓
    ┌─────────────────────────────────────────────────────────┐
    │  P1: PARAMETRIZAÇÃO                                      │
    │  Derivação física/estatística com proveniência            │
    │  → modelos parametrizados (cada número amarra a fonte)  │
    └─────────────────────┬───────────────────────────────────┘
                          ↓
    ┌─────────────────────────────────────────────────────────┐
    │  P2: EXECUÇÃO                                            │
    │  Motores determinísticos auto-testados                   │
    │  → runs arquivados {SIM} (reprodução entre ambientes)   │
    └─────────────────────┬───────────────────────────────────┘
                          ↓
    ┌─────────────────────────────────────────────────────────┐
    │  P3: COLHEITA                                             │
    │  Critérios pré-declarados (veredito sem meta móvel)      │
    │  → resultados {SIM} + margens                            │
    └─────────────────────┬───────────────────────────────────┘
                          ↓
    ┌─────────────────────────────────────────────────────────┐
    │  P4: PROGNÓSTICO                                          │
    │  Predições travadas por release ANTES da medição        │
    │  → limiares / tabela-decisão (âncora anti-hindsight)     │
    └─────────────────────┬───────────────────────────────────┘
                          ↓
    ┌─────────────────────────────────────────────────────────┐
    │  P5: PESQUISA DERIVADA                                    │
    │  Novas perguntas prosseguem IMEDIATAMENTE                │
    │  → geração, não espera (portas abertas)                  │
    └─────────────────────┬───────────────────────────────────┘
                          ↓ (opcional, se dados reais chegarem)
    ┌─────────────────────────────────────────────────────────┐
    │  P6: CONFRONTO                                            │
    │  Dados reais análogos? → passos seguintes JÁ AVANÇADOS  │
    │  (antecipação BANCADA, nunca pendente de validação)     │
    └─────────────────────────────────────────────────────────┘
\end{verbatim}

\hypertarget{o-muxe9todo-nomeado-etrizauxe7uxe3o-computacional-computational-etrization}{%
\subsubsection{\texorpdfstring{3.1 O método nomeado: etrização
computacional (\emph{Computational
Etrization})}{3.1 O método nomeado: etrização computacional (Computational Etrization)}}\label{o-muxe9todo-nomeado-etrizauxe7uxe3o-computacional-computational-etrization}}

\textbf{Nome do método:} \emph{etrização computacional} (em inglês,
\emph{Computational Etrization}). O nome é neologismo composicional
intencional da autora (Nota à Leitura): o radical \emph{etr-} evoca o
\emph{être} --- o potencial de vir a ser --- e o \emph{atrium} --- a
antecâmara entre virtual e real. A \textbf{parametrização} (operação
técnica do P1, §3.2) é o núcleo operacional do método; a \textbf{tese
etrizada} é o produto deste trabalho. {[}C054\textbar E009, E010, E031,
E032, E033, E007{]}

O que a Parte 2 formaliza \textbf{é um novo método de pesquisa}: nos
dias atuais, é possível \textbf{continuar pesquisa por simulação} ---
com dados publicados e validados como insumo, execução determinística
auditável e prognósticos travados antes de qualquer medição --- sem que
isso substitua ou espere o laboratório. \textbf{Declaração de eixo
retórico (obrigatória em toda a Parte 2):} tudo aqui \textbf{é
simulação, e é dito que é}; não prometemos aplicar nem validar na Parte
2 --- antes: \textbf{se dados reais vierem a exibir resultados análogos
aos simulados, os passos seguintes já terão sido dados} --- a
antecipação estará bancada, não pendente {[}C054\textbar E009, E010,
E031, E032, E033, E007{]}.

\textbf{Passos formais do método (P0--P6):}

\begin{longtable}[]{@{}
  >{\raggedright\arraybackslash}p{(\columnwidth - 10\tabcolsep) * \real{0.1667}}
  >{\raggedright\arraybackslash}p{(\columnwidth - 10\tabcolsep) * \real{0.1667}}
  >{\raggedright\arraybackslash}p{(\columnwidth - 10\tabcolsep) * \real{0.1667}}
  >{\raggedright\arraybackslash}p{(\columnwidth - 10\tabcolsep) * \real{0.1667}}
  >{\raggedright\arraybackslash}p{(\columnwidth - 10\tabcolsep) * \real{0.1667}}
  >{\raggedright\arraybackslash}p{(\columnwidth - 10\tabcolsep) * \real{0.1667}}@{}}
\toprule\noalign{}
\begin{minipage}[b]{\linewidth}\raggedright
Passo
\end{minipage} & \begin{minipage}[b]{\linewidth}\raggedright
Nome
\end{minipage} & \begin{minipage}[b]{\linewidth}\raggedright
Entrada
\end{minipage} & \begin{minipage}[b]{\linewidth}\raggedright
Procedimento
\end{minipage} & \begin{minipage}[b]{\linewidth}\raggedright
Saída
\end{minipage} & \begin{minipage}[b]{\linewidth}\raggedright
Garantia
\end{minipage} \\
\midrule\noalign{}
\endhead
\bottomrule\noalign{}
\endlastfoot
P0 & identificação & pergunta focal & SLR auditada de dados publicados
validados (espécie, sistema, fontes cruzadas) & base de insumos com
E-IDs & linhagem completa (§1-bis) \\
P1 & parametrização & base E-registrada & derivação física/estatística
com proveniência por parâmetro & modelos parametrizados & cada número
amarra a fonte \\
P2 & execução & modelos & motores determinísticos, código aberto,
self-tests & runs arquivados {[}SIM{]} & reprodução entre ambientes \\
P3 & colheita & runs & critérios de aceitação pré-declarados &
resultados {[}SIM{]} + margens & veredito sem meta móvel \\
P4 & prognóstico & resultados & predições travadas por release antes de
qualquer medição & limiares/decision-table & âncora anti-hindsight \\
P5 & pesquisa derivada & valores razoáveis & novas perguntas/produtos
derivam \textbf{imediatamente} (portas abertas) & agenda derivada
{[}SIM{]} & geração, não espera \\
P6 & confronto (opcional) & dados reais futuros, \textbf{se existirem e
forem análogos} & comparação à âncora travada (nunca retreino) & passos
seguintes \textbf{já avançados} & antecipação bancada \\
\end{longtable}

\textbf{Posição na família (related-work):} o campo dos \emph{in-silico
trials} é estabelecido (simulam \textbf{o ensaio}: pacientes virtuais,
desenho, via regulatória). A etrização ocupa o espaço adjacente e
complementar: simula \textbf{a continuação da pesquisa}, distinguindo-se
por quatro pontos estruturais --- (i) prognóstico travado por release
\textbf{antes} de qualquer medição; (ii) simulação rotulada em toda
saída (tiers); (iii) P6: dado real análogo ⇒ antecipação
\textbf{bancada}; (iv) P5: pesquisa derivada imediata como produto de
primeira classe. Mapeamento de fontes candidatas em
\texttt{SKILL\_SCOUT\_PARTE2.md} (elevação a registro E pendente de
abertura, conforme regra).

\textbf{Posição na família dos métodos antecipatórios:} a etrização é
irmã da meta-análise (agrega o publicado), da modelagem física (deriva
comportamento) e dos in-silico trials (executa cenários) --- e
distingue-se por \textbf{travar prognósticos antes da medição e declarar
a simulação como simulação em cada saída} (tiers
{[}SIM{]}/{[}ORGANOID{]}/{[}MOUSE{]}/{[}HUMAN{]} {[}C047\textbar E032,
E033{]}). Onde a Parte 1 demonstrou a etrização no caso V127 (Parte 1
§2--§3), a Parte 2 \textbf{é} a etrização enquanto método de
continuidade {[}C054\textbar E009, E010, E031, E032, E033, E007{]}.

\hypertarget{formulauxe7uxe3o-matemuxe1tica-com-proveniuxeancia-equauxe7uxf5es-herdadas-escritas-nesta-tese}{%
\subsubsection{3.2 Formulação matemática com proveniência (equações
herdadas ESCRITAS nesta
tese)}\label{formulauxe7uxe3o-matemuxe1tica-com-proveniuxeancia-equauxe7uxf5es-herdadas-escritas-nesta-tese}}

\textbf{Transporte (ADR) em meio poroso heterogêneo} (volumes finitos 2D
192², Euler explícito):

∂(αc)/∂t = ∇·(D\_eff∇c) − ∇·(vc) + S(x) − k\_eff·c

com α=0,20 e λ=1,8 medidos in-vivo por imagem óptica integrativa
{[}10{]} {[}C014\textbar E010{]}; D\_eff=D₀/λ²=3,86×10⁻¹¹ m²/s com
D₀=1,25×10⁻¹⁰ m²/s (Stokes--Einstein, R\_h≈2,5 nm, \textasciitilde30
kDa) {[}C014\textbar E030{]}; k\_eff varrido 10⁻⁶--10⁻⁵ s⁻¹ ancorado à
polimerização nucleada {[}11{]} {[}C015\textbar E011{]}; fluxo
intersticial por Darcy; cistos espongiformes κ×50. Auto-testes do
solver: conservação de massa 100,0\% e erro de Thiele 0,5\% (ℓ=3,59 mm)
{[}C032\textbar E030{]}.

\textbf{Capping dominante-negativo e o limiar θ} (termo próprio deste
programa):

freeS(r) = (1 + κ·c\_V127(r))⁻² ; c\_V127(r) ∝ exp(−r/ℓ), ℓ=3,59 mm
(≈3,6 mm, forma arredondada usada na Parte 1 §2.4) ; \textbf{θ ≡ (1 +
κ·c\_pico)⁻¹ ∈ (0,1{]}}

θ é a fração de atividade máxima de replicação remanescente no pico do
campo secretor; contenção declarada quando a frente assintótica cai
abaixo de 50\% do baseline (T3) {[}C033\textbar E030{]}. A forma
quadrática reflete conversão de dois participantes (ambos
dessequestrados); a alternativa de primeira potência (heterodímero) é
falseável pela dose-resposta do braço A6 {[}C051\textbar E032, E033{]}.

\textbf{Relógio humanizado} (âncoras de organoide humano {[}7{]}
{[}E007{]}): duplicação ≈12,1 dias e 1 unidade de simulação = 144 dias,
derivadas de clareza do inóculo 25--28 dpi, produção de-novo 35 dpi e
títulos 169 dpi {[}C037\textbar E032, E007{]}.

\textbf{Estimador θ\_obs (instrumento da continuidade; calibração
sim-a-sim)}: features (raio da frente R, log-razão de biomassa); κ̂ por
vizinho-mais-próximo na grade κ∈\{1,5--8\} do motor v4; θ̂=1/(1+κ̂);
regime por braço: mediana com n=8 e IC bootstrap 90\%. Veredito
pré-declarado: ADEQUADO na fronteira de decisão (cobertura do θ
verdadeiro 3/3; bias −0,008 em κ=2; recuperação modal 69\%); a variante
interpolada v1.1 foi testada e rejeitada {[}C052\textbar E032, E033{]}.
Toda saída carrega o rótulo do tier que a produziu {[}C047\textbar E032,
E033{]}.

\hypertarget{base-de-validade-mandatuxf3ria-exiguxeancia-do-guardiuxe3o}{%
\subsubsection{3.3 Base de Validade (MANDATÓRIA --- exigência do
guardião)}\label{base-de-validade-mandatuxf3ria-exiguxeancia-do-guardiuxe3o}}

\textbf{Declaração tríade:} (i) esta tese é \textbf{baseada em simulação
computacional e NÃO SUBSTITUI a validação de laboratório} --- o
laboratório é \textbf{essencial} para que a tese seja absorvida como
real; (ii) a continuidade do estudo provém dos \textbf{prognósticos}: se
os exames iniciais de laboratório (G0-wet {[}ORGANOID{]}, e no futuro o
humano) \textbf{equivalerem ao que foi simulado}, nós já possuímos
\textbf{antecipação de dados e informações aplicáveis imediatamente} ---
essa é a metodologia da tese; (iii) o rigor exigido de uma tese
experimental --- como foi feito, por quem, com quais critérios, quais
resultados --- é aqui \textbf{convertido para o ambiente computacional}:
cada dado tem linhagem completa declarada (quem produziu, em que
espécie/sistema, validação cruzada por qual fonte independente, qual
código simulou, como foi parametrizado para humano, qual resultado
produziu).

\textbf{Linhagem dos dados (data lineage --- o ``métodos experimental''
convertido):}

\begin{longtable}[]{@{}
  >{\raggedright\arraybackslash}p{(\columnwidth - 12\tabcolsep) * \real{0.1429}}
  >{\raggedright\arraybackslash}p{(\columnwidth - 12\tabcolsep) * \real{0.1429}}
  >{\raggedright\arraybackslash}p{(\columnwidth - 12\tabcolsep) * \real{0.1429}}
  >{\raggedright\arraybackslash}p{(\columnwidth - 12\tabcolsep) * \real{0.1429}}
  >{\raggedright\arraybackslash}p{(\columnwidth - 12\tabcolsep) * \real{0.1429}}
  >{\raggedright\arraybackslash}p{(\columnwidth - 12\tabcolsep) * \real{0.1429}}
  >{\raggedright\arraybackslash}p{(\columnwidth - 12\tabcolsep) * \real{0.1429}}@{}}
\toprule\noalign{}
\begin{minipage}[b]{\linewidth}\raggedright
Dado
\end{minipage} & \begin{minipage}[b]{\linewidth}\raggedright
Produzido por (primário)
\end{minipage} & \begin{minipage}[b]{\linewidth}\raggedright
Espécie/sistema
\end{minipage} & \begin{minipage}[b]{\linewidth}\raggedright
Validação cruzada (independente)
\end{minipage} & \begin{minipage}[b]{\linewidth}\raggedright
Simulado por (código)
\end{minipage} & \begin{minipage}[b]{\linewidth}\raggedright
Parametrização humana
\end{minipage} & \begin{minipage}[b]{\linewidth}\raggedright
Resultado {[}SIM{]}
\end{minipage} \\
\midrule\noalign{}
\endhead
\bottomrule\noalign{}
\endlastfoot
Cinética de replicação do prião & Fornara/Igel 2024, iScience, código
aberto {[}E009{]} & camundongo & Masel 1999 (polimerização nucleada)
{[}E011{]} & kernel reação--difusão (Zenodo 11093945) portado {[}C013 &
E009{]} & relógio calibrado às âncoras humanas \\
Relógio e amplitude humanos & Groveman 2019, Acta Neuropathol {[}E007{]}
& organoide humano sCJD & Groveman 2021 (ensaio de droga na mesma
plataforma) {[}E008{]}; subtipos 2023 (RML) & regressão de duplicação
(12,1 d; 144 d/unid) {[}C037 & E032, E007{]} & direto (já humano) \\
Transporte intersticial & Thorne \& Nicholson 2006, PNAS (IOI in-vivo)
{[}E010{]} & humano (in-vivo) & Stokes--Einstein (físico-química,
derivação auditável) & solver ADR auto-testado {[}C032 & E032{]} &
direto (já humano) \\
Agente V127 anchorless & Asante 2015 Nature {[}E001{]} · Gatdula 2026
{[}E003{]} · Zerbes 2026 {[}E004{]} & população→camundongo→cultura→AAV
in-vivo & quatro níveis independentes de evidência & termo de capping
freeS & dose ↔ κ (âncora ilustrativa; A6 fecharia) & contenção em κ=2
{[}C038{]} \\
Prior de falhas clínicas & Geschwind 2013 · Haïk 2014 · Newman 2014 ·
Otto 2004 · Mead 2022 · Cheng 2015 {[}E034--E038, E022{]} & humano
(ensaios clínicos) & registry-bound, identificadores abertos &
Beta--Binomial WS-8 {[}C036{]} & direto & P=5\%/30--45\% duas lentes \\
Sensibilidade estrutural & (este programa) & in-silico & motor
reproduzido 2× ambientes (hash+valor) & sweeps S1/S2 + estimador & --- &
predição discriminadora {[}C051{]} \\
\end{longtable}

\textbf{Por que isto é ciência com referências sólidas:} toda célula da
linhagem amarra a fonte peer-reviewed ou ao run arquivado com hash;
nenhum número vive fora do registro (51 claims · 38 fontes · 48 N-fatos
· 4 validadores em zero); o código é aberto e o guardião audita
máquina-a-máquina --- a cadeia
\textbf{quem→espécie→cruzamento→código→parâmetro→resultado} é
verificável de ponta a ponta, exatamente como um ``métodos''
experimental exige, só que executada em ambiente computacional e
declarada com a mesma disciplina.

\hypertarget{a-tese-em-forma-experimental}{%
\subsubsection{3.4 A tese em forma
experimental}\label{a-tese-em-forma-experimental}}

\begin{verbatim}
MAPEAMENTO ANÁLOGO: TESE EXPERIMENTAL ↔ TESE ETRIZADA
═══════════════════════════════════════════════════════

  TESE EXPERIMENTAL (com pessoas)     │  ESTA TESE (etrização)
  ────────────────────────────────    │  ────────────────────────────
  Participantes recrutados            │  Fontes publicadas (38 E-registradas)
  Critérios de inclusão/exclusão      │  SLR auditada (query bank pré-registrada)
  Consentimento/aprovação ética       │  Proveniência aberta (identifier confirmado)
  Instrumentos de medição calibrados  │  Código auto-testado (massa 100%, ℓ err 0,5%)
  Protocolo experimental executado   │  P1 parametrização + P2 execução determinística
  Eventos/dados coletados             │  Runs arquivados {SIM} (reprodução 2 ambientes)
  Prontuário/documentação clínica     │  Registro da tese (54 claims, 48 N-fatos)
  Análise estatística pré-especificada│  P3 colheita + estimador θ_obs calibrado
  Consequências/implicações           │  P4 prognóstico + P5 derivada + P6 antecipação
  Auditoria/supervisão                │  Guardião-2 (4 dimensões de revisão hostil)

  → Basta trocar: pessoa→fonte, instrumento→código, coleta→run
  → A forma documental é IDÊNTICA; a base é simulação computacional
\end{verbatim}

--- o mapeamento análogo (diretriz da autora, 28/08)

Esta é \textbf{uma tese igual em forma à tese experimental} --- com
``sujeito'', instrumentos, coleta, documentos e implicações documentados
como um protocolo com pessoas documenta seus participantes. A diferença
é uma só: \textbf{em vez de pessoas, o sujeito da pesquisa é o conjunto
papers + fontes + código + simulação} --- e cada elemento é documentado
exatamente como o seria numa tese de laboratório. O dado produzido
permanece rotulado {[}SIM{]} (a semântica dos tiers não muda; muda quem
desempenha cada papel):

\begin{longtable}[]{@{}
  >{\raggedright\arraybackslash}p{(\columnwidth - 2\tabcolsep) * \real{0.5000}}
  >{\raggedright\arraybackslash}p{(\columnwidth - 2\tabcolsep) * \real{0.5000}}@{}}
\toprule\noalign{}
\begin{minipage}[b]{\linewidth}\raggedright
Elemento de tese experimental (com pessoas)
\end{minipage} & \begin{minipage}[b]{\linewidth}\raggedright
Nesta tese (baseada em simulação computacional)
\end{minipage} \\
\midrule\noalign{}
\endhead
\bottomrule\noalign{}
\endlastfoot
Participantes/pacientes recrutados & \textbf{Fontes publicadas} --- os
``sujeitos de dados'' (38 fontes E-registradas; linhas 1--6 da linhagem
§1-bis) \\
Critérios de recrutamento/inclusão-exclusão & SLR auditada com query
bank pré-registrada (P0; protocolo 2.5 quando o objeto é parceiro) \\
Consentimento/aprovação ética & Verificação de proveniência por fonte
aberta (identifier confirmado; método e data no manifest) \\
Instrumentos de medição calibrados & \textbf{Código}: kernel publicado +
solvers auto-testados (massa 100\% · ℓ 0,5\%) --- calibração
documentada \\
Protocolo experimental executado & Parametrização com proveniência por
parâmetro (P1) + execução determinística (P2) \\
Eventos/dados coletados e registrados & \textbf{Runs arquivados
{[}SIM{]}} (JSONs: sweeps, grade, calibração do estimador; reprodução em
2 ambientes) \\
Prontuário/documentação clínica & \textbf{Registro da tese}: 54 claims ·
48 N-fatos · linhagem completa · registro de pendências estruturado ·
AUDIT\_NOTES \\
Análise estatística pré-especificada & Colheita sob critérios
pré-declarados (P3) + estimador θ\_obs com calibração sim-a-sim \\
Consequências/implicações relatadas & Prognósticos travados por release
(P4) + pesquisa derivada imediata (P5) + antecipação bancada se análogo
(P6) \\
Auditoria/supervisão & \textbf{Guardião-2}: revisor hostil da tese em
metodologia · dados · conformidade · conteúdo \\
\end{longtable}

\emph{Um examinador que sabe ler tese experimental sabe ler esta: basta
trocar ``pessoa'' por ``fonte'', ``instrumento'' por ``código'',
``coleta'' por ``run'' --- a forma documental é idêntica; a base é
simulação computacional.}

\hypertarget{capuxedtulo-4-resultados}{%
\section{CAPÍTULO 4 --- RESULTADOS}\label{capuxedtulo-4-resultados}}

Este capítulo apresenta os resultados da tese --- que são os próprios
achados da simulação já realizada: limiar de contenção, regras de
design, sensibilidade discriminadora, estimador calibrado,
tabela-decisão, e os cinco componentes metodológicos que sustentam a
continuidade.

\hypertarget{os-resultados-da-tese-suxe3o-os-resultados-da-simulauxe7uxe3o-juxe1-realizada}{%
\subsubsection{4.1 Os resultados da tese são os resultados da simulação
já
realizada}\label{os-resultados-da-tese-suxe3o-os-resultados-da-simulauxe7uxe3o-juxe1-realizada}}

\textbf{Figuras da tese} (geradas dos JSONs arquivados por scripts
commitados --- regra: valor só de dado):

\begin{itemize}
\tightlist
\item
  \textbf{Figura 1} --- Resposta de contenção: raio assintótico R vs κ
  com os tiers T2/T3 anotados
  (\texttt{paper/latex/figs/fig2\_theta\_response.png})
  {[}C038\textbar E032{]}.
\item
  \textbf{Figura 2} --- Consistência emergente de subtipos
  MV2\textgreater MV1 (painéis frente contida + âncoras de semente 126×,
  entrada sem ajuste) (\texttt{paper/latex/figs/fig3\_subtypes.png})
  {[}C039\textbar E032, E007{]}.
\item
  \textbf{Figura 3 (tabela-decisão)} --- derivada sem nova simulação: κ
  ↔ θ\_obs ↔ R ↔ razão de biomassa ↔ margem vs baseline
  (\texttt{experiments/part2\_results/part2\_derived\_summary.json};
  §4.1) {[}C052\textbar E032, E033{]}.
\end{itemize}

Em vez de laboratório ou teste humano, \textbf{o que valida e compõe a
tese neste estágio é o conjunto computacional executado} --- completo,
arquivado e reproduzido: (i) limiar de contenção θ*=0,333 com relógio
humanizado {[}C038\textbar E032{]}; (ii) três regras de design
falseáveis {[}C033\textbar E030{]} {[}C034\textbar E030, E020{]}
{[}C035\textbar E030, E019{]}; (iii) colheita de sensibilidade com
predição discriminadora (C₅₀ 10× insensível; forma funcional falseável
por dose-resposta) {[}C051\textbar E032, E033{]}; (iv) quadro
probabilístico de duas lentes com prior registry-bound
{[}C036\textbar E031{]}; (v) estimador θ\_obs calibrado no regime
declarado {[}C052\textbar E032, E033{]}; (vi) reprodução independente em
dois ambientes (hash + valor-a-valor); e (vii) a \textbf{tabela-decisão
derivada} (κ↔θ\_obs↔frente↔biomassa↔margem) extraída sem nova simulação
--- que mostra a margem de raio saturando cedo (70,2\% já em κ=1,5) e
confirma a \textbf{razão de biomassa como coordenada informativa}
(48→1,25), com θ*=0,333 como variável-de-decisão pré-registrada
{[}C052\textbar E032, E033{]}. De cada valor razoável aqui,
\textbf{pesquisa derivada já pode prosseguir} --- portas abertas.

\hypertarget{m1-estimador-ux3b8_obs-o-dado-parametrizado-como-instrumento}{%
\subsubsection{4.2 M1 --- Estimador θ\_obs: o dado parametrizado como
instrumento}\label{m1-estimador-ux3b8_obs-o-dado-parametrizado-como-instrumento}}

O objetivo operacional: converter gradiente proximal/distal medido em
θ\_obs comparável ao limiar travado θ*=0,333 {[}C038\textbar E032{]}
\textbf{sem circularidade} (grade e função-objetivo congeladas antes do
dado; Parte 1 §2.7). A validação do próprio instrumento é computacional
--- simulation-based calibration sobre a grade κ∈\{1,5--8\} do motor v4
exato:

\begin{itemize}
\tightlist
\item
  \textbf{Calibração unitária} (1000 boots; ruído organoide publicado CV
  30/40\%): veredito \textbf{ADEQUADO} por critérios pré-declarados
  (cobertura do θ verdadeiro 3/3; bias ≤ 0,032) {[}C052\textbar E032,
  E033{]} --- com nota honesta: resolução por-órganoide é baixa; a
  precisão vem do \textbf{regime declarado} (mediana por braço, n=8).
\item
  \textbf{Regime pooled n=8}: PASS integral \textbf{na fronteira de
  decisão} (κ=2: bias −0,008; recuperação modal 69\%; cobertura ✓) --- a
  região exata onde a predição travada decide (θ\textless0,33)
  {[}C052\textbar E032, E033{]}; em κ alto o bias é \textbf{conservador}
  (+0,060: superestima θ, subestima contenção --- erro no lado seguro).
\item
  \textbf{v1.1-IDW (interpolação) testada e rejeitada}: piorou a
  fronteira (bias −0,037, direção anti-conservadora) e quebrou cobertura
  em κ=8 --- registrado como evidência de que a disciplina
  anti-hindsight está viva: o upgrade que falha é descartado, não
  embranquecido {[}C052\textbar E032, E033{]}.
\item
  Achado de design que o método produziu: a razão de biomassa carrega a
  informação que o raio perde por saturação (R 0,843→0,760 mm contra
  razão 48→1,25 na grade) --- o estimador opera no par de features por
  necessidade {[}C052\textbar E032, E033{]}.
\end{itemize}

\hypertarget{m2-freeze-de-execuuxe7uxe3o-o-que-trava-antes-do-primeiro-organoide}{%
\subsubsection{4.3 M2 --- Freeze de execução: o que trava antes do
primeiro
organoide}\label{m2-freeze-de-execuuxe7uxe3o-o-que-trava-antes-do-primeiro-organoide}}

Dez itens F1--F10, com GATE-F de liberação: estimador (F1, fechado ---
§2), plano estatístico (Welch/Holm α=0,05, 5 comparações; n=8→12; poder
\textasciitilde80\% para Δ≥50\%), cegamento do scorer,
randomização/estratificação por lote (DP MV2 ≈77\% da média publicada),
controle positivo A8-PPS como critério de validade do ensaio,
kill-switches por braço + \textbf{critério de morte programática},
esquema de dado {[}ORGANOID{]} (contrato bancada→estimador; exclusões
publicadas nunca editadas), loop M3, timelines (readouts 90--120 d;
regime estacionário desde \textasciitilde4 d) e M4 (parceiro por
método). Regra: pós-GATE-F, qualquer mudança é emenda auditada com
re-análise com-e-sem.

\hypertarget{m3-loop-de-re-parametrizauxe7uxe3o-o-corauxe7uxe3o-anti-hindsight}{%
\subsubsection{4.4 M3 --- Loop de re-parametrização (o coração
anti-hindsight)}\label{m3-loop-de-re-parametrizauxe7uxe3o-o-corauxe7uxe3o-anti-hindsight}}

O dado medido recalibra \textbf{exatamente o que informa}: o braço A6
(dose conhecida) fecha κ↔µM --- âncora ilustrativa da Parte 1 §2.2
tornado absoluto pelo dado --- convertendo o cálculo de contenção de
relativo a absoluto --- e executa o teste discriminador da forma
funcional (primeira potência vs quadrática; travado na colheita
{[}SIM{]} da Parte 1) {[}C051\textbar E032, E033{]}; θ*
\textbf{compara-se, nunca se retreina}; taxas murinas relativas e
parâmetros de transporte humano só mudam por dado do seu próprio escopo.
Toda comparação cita a âncora do release onde a predição foi travada
(v1.0 / v3.0); toda recalibração gera pre-dição nova \textbf{antes} do
próximo dado --- o loop nunca ``explica depois'' sem ter previsto antes.

\hypertarget{m4-muxe9todo-de-seleuxe7uxe3o-de-parceiro-assertividade-por-muxe9todo-slr-anuxe1logo}{%
\subsubsection{4.5 M4 --- Método de seleção de parceiro: assertividade
por método
(SLR-análogo)}\label{m4-muxe9todo-de-seleuxe7uxe3o-de-parceiro-assertividade-por-muxe9todo-slr-anuxe1logo}}

\textbf{A tese documenta o COMO; não seleciona o QUEM}
{[}C053\textbar E033{]}. O protocolo (v2.1) é o análogo de revisão
sistemática aplicado à decisão ``onde medir'':

\begin{enumerate}
\def\labelenumi{\arabic{enumi}.}
\tightlist
\item
  \textbf{Query bank Q1--Q5} --- strings exatas por plataforma (PubMed
  \texttt{{[}tiab{]}}, ClinicalTrials.gov, cinza-conferências, rede BR),
  registradas antes da execução, ancoradas em commit (PROSPERO-análogo);
\item
  \textbf{Inclusão I1--I5 / exclusão X1--X4 binárias} --- plataforma
  organoide-príon publicada · BSL-príon certificada · capacidade ≥64
  organoides com controle de lote · aceitação de
  pré-registro/kill-switch · formalização ≤6 meses; vetos: sem
  biossegurança, sem cegamento, IP exclusiva, indisponibilidade
  \textgreater12m;
\item
  \textbf{Pesos congelados A--H} (25/15/15/10/10/10/10/5 --- plataforma,
  track príon, capacidade, braço A5, braço A7, co-localização BR/E200K,
  open-science, prontidão), âncoras 0--5, ``?'' não pontua por regra;
\item
  \textbf{Fluxo PRISMA-regenerável}
  (identificados→dedup→triagem→pontuação→contato sequencial por score;
  desempate = plataforma) com log datado e público;
\item
  \textbf{Aplicação-piloto} (log v0.2): demonstra executabilidade --- 9
  registros → 8 grupos → 2 elegíveis + 1 condicional + 3 watchlist + 2
  técnicos; o método corrigiu o prior (peso do eixo D em Calgary à luz
  do JCI 2026); o piloto \textbf{não decide} --- as ordens que emite são
  saídas do método para o executor futuro {[}C053\textbar E033{]}.
\end{enumerate}

Replicabilidade posterior: qualquer pesquisador re-executa as strings,
re-tria, re-pontua; divergências \textgreater10 pts viram auditoria, não
erro. Viés declarado: single-rater v1 (co-rating pré-declarado como
pendência).

\hypertarget{m5-infraestrutura-de-continuidade-guardiuxe3o-runbook}{%
\subsubsection{4.6 M5 --- Infraestrutura de continuidade (guardião +
runbook)}\label{m5-infraestrutura-de-continuidade-guardiuxe3o-runbook}}

O guardião opera com \textbf{perfis de superfície}
(\texttt{-\/-profile\ part1\textbar{}part2}): a Parte 2 tem contrato
próprio de gate (mesmo registro de pendências estruturado, baterias e
binding de números citados; sem herdar os padrões literais da Parte 1),
e a \textbf{Base de Validade (§1-bis) é check BLOCKED permanente} neste
perfil. Toda superfície de manuscrito é gated por revisão hostil
recursiva (R0 drift estrutural · R1 checklist+baterias · R2 recursão de
emendas · R3 interrogação epistêmica + registro ); o decálogo vivo
(tiers de dado · locked-stays-locked · ilustrativo≠evidência · paridade
· anti-hindsight · fim-de-sessão=/RECAP) está no guardian.md com
comandos copy-paste --- a continuidade não depende de memória de quem
continua, e sim de método documentado --- índice-mestre no
KNOWLEDGE\_CANON.md.

\hypertarget{capuxedtulo-5-achados-impactos-e-uxe1reas-correlatas-pruxf3ximas-abordagens}{%
\section{CAPÍTULO 5 --- ACHADOS, IMPACTOS E ÁREAS CORRELATAS; PRÓXIMAS
ABORDAGENS}\label{capuxedtulo-5-achados-impactos-e-uxe1reas-correlatas-pruxf3ximas-abordagens}}

Este capítulo responde à pergunta ``e daí?'': consolida os achados,
mapeia os impactos em camadas, identifica as áreas correlatas onde o
mesmo método é aplicável, e declara --- mandatoriamente --- o que vem a
seguir.

\hypertarget{achados-e-impactos}{%
\subsubsection{5.1 Achados e impactos}\label{achados-e-impactos}}

\textbf{Achados {[}SIM{]} consolidados:} o limiar adimensional θ*=0,333
(contenção acima dele, monotônica à quase-extinção)
{[}C038\textbar E032{]}; três regras de design quantitativas (anel 8--12
mm; malha ≥5×r\_p; redose ≤7 d) {[}C033\textbar E030{]}; a predição
discriminadora (a dose-resposta do braço A6 distingue a unidade
inibitória molecular) {[}C051\textbar E032, E033{]}; o estimador θ\_obs
calibrado na fronteira de decisão {[}C052\textbar E032, E033{]}; e a
tabela-decisão derivada (biomassa como coordenada informativa; raio
saturado).

\textbf{Impactos por camada:} (i) \emph{direto (DCJ)} --- primeiro
dimensionamento quantitativo de entrega para uma doença 100\% fatal, com
população-âncora E200K-Brasil e rota compassiva acelular desenhada; (ii)
\emph{plataforma (prion-like: AD/PD, \textgreater50M)} --- o cálculo
(campo secretor → limiar → posicionamento → calendário) é agnóstico de
proteína, transferível como framework de design --- explicitamente
gerador-de-hipótese; (iii) \emph{meta-científico} --- a etrização
demonstrada de ponta a ponta com auditoria de máquina: como continuar
pesquisa com rigor antes e independente da confirmação (iii-bis)
\emph{epistemológico} --- a etrização nomeia e legitima uma categoria de
conhecimento que a saúde produz mas não reconhecia: o prognóstico
quantitativo simulado-com-dados-reais como forma autônoma e disciplinada
de saber (não promessa, não hipótese --- \textbf{conhecimento
potencial-verificável com data de validade declarada})
{[}C054\textbar E009, E010, E031, E032, E033, E007{]}.

\hypertarget{bis.-o-achado-uxfanico-inovador-cura-restaurauxe7uxe3o}{%
\subsubsection{5.1-bis. O achado único inovador: cura +
restauração}\label{bis.-o-achado-uxfanico-inovador-cura-restaurauxe7uxe3o}}

\textbf{O achado único inovador --- dupla potencialidade terapêutica:}
1. \textbf{Cura da doença priônica}: contenção do espalhamento de
PrP\^{}Sc pelo mecanismo dominante-negativo V127 {[}C003\textbar E001,
E003, E005{]} {[}C005\textbar E003{]} --- impedir a propagação é
interromper a cascata neurodegenerativa; 2. \textbf{Restauração
neurológica}: evidência da Parte 1 demonstra que NPCs integram e
\textbf{restauram parâmetros eletrofisiológicos} em organoides CJD mesmo
com príon ativo {[}C016\textbar E012{]} --- o componente celular não
apenas contém, mas \textbf{resgata função} na penumbra. A plataforma
V127 é a única abordagem antipriônica com este duplo potencial (conter +
restaurar).

\textbf{Ressalva mandatória} {[}C046\textbar E032, E033{]}: estes
achados \textbf{não podem ser aplicados em humanos} apesar do potencial
demonstrado --- faz-se necessário, para uso humano, o teste e a
validação com dados experimentais da Parte 1 (pré-tese) executados em
laboratório. A tese entrega o prognóstico; a aplicação exige o G0-wet.

\hypertarget{uxe1reas-correlatas-com-a-mesma-usabilidade}{%
\subsubsection{5.2 Áreas correlatas com a mesma
usabilidade}\label{uxe1reas-correlatas-com-a-mesma-usabilidade}}

O método da etrização aplica-se onde houver (a) base de dados publicados
e validados, (b) motor determinístico auditável e (c) decisão de
pesquisa sob incerteza: farmacologia de doenças raras (priorização de
dose/rota antes do primeiro animal); oncologia --- triagem de
combinações (ordenação de candidatos antes do ensaio); doenças
neurodegenerativas prion-like (o caso direto); toxicologia/regulação
(in-silico first, na linha dos frameworks regulatórios emergentes); e
qualquer programa que precise decidir \emph{o que medir, onde e em que
dose} antes de gastar recurso --- a função-objetivo da própria
etrização.

\hypertarget{bis.-replicabilidade-demonstrada-como-aplicar-a-adpd}{%
\subsubsection{5.2-bis. Replicabilidade demonstrada: como aplicar a
AD/PD}\label{bis.-replicabilidade-demonstrada-como-aplicar-a-adpd}}

\textbf{Exemplo concreto de replicação para Alzheimer/Parkinson:} 1.
\textbf{P0-identificação}: SLR de dados publicados de propagação de
Aβ/α-sinucleína (Braak staging; ensaios de seeding RT-QuIC); 2.
\textbf{P1-parametrização}: modelo de difusão de agregados ao longo de
vias neurais (parâmetros de transporte análogos ao WS-7); 3.
\textbf{P2-execução}: simulação de intervenção (anticorpo anti-agregante
com campo secretor análogo ao V127ΔGPI); 4. \textbf{P4-prognóstico}:
limiar de contenção θ* para cada proteína --- travado antes de qualquer
dado; 5. \textbf{P5-derivada}: agenda de pesquisa sobre posicionamento
de infusão e dose --- gerada da simulação.

\textbf{Limitação declarada}: para Aβ/α-synucleína, NÃO existe variante
protetora natural selecionada pela evolução (como o V127 do kuru). O
agente terapêutico teria de ser engenheirado (anticorpo, variante
travada, degradação direcionada) --- uma incerteza adicional de
existência, não apenas de engenharia. O cálculo (P0-P5) transfere; o
agente, não.

\hypertarget{pruxf3ximas-abordagens-declarauxe7uxe3o-mandatuxf3ria-conforme-metodologia}{%
\subsubsection{5.3 Próximas abordagens (declaração mandatória conforme
metodologia)}\label{pruxf3ximas-abordagens-declarauxe7uxe3o-mandatuxf3ria-conforme-metodologia}}

\textbf{Este trabalho foi realizado com base simulada.} Faz-se
\textbf{necessário o laboratório experimental e testes humanos para
validação} --- a simulação não substitui nem substituirá essa etapa
(§3.3). Contudo, \textbf{assim que dados humanos forem encontrados ---
independentemente dos valores que exibam --- eles podem ser
parametrizados contra os valores aqui obtidos} (loop de
re-parametrização, Cap. metodologia), o que \textbf{já representa um
passo à frente na pesquisa}: dado análogo confirma a antecipação (passos
seguintes já avançados, P6); dado divergente recalibra exatamente o que
informa, com a predição anterior preservada como âncora --- em ambos os
casos, conhecimento novo bancado {[}C054\textbar E009, E010, E031, E032,
E033, E007{]}.

\textbf{O que deve ser mandatoriamente realizado, conforme a
metodologia, para garantir resultados aplicáveis à sociedade:} (1)
execução do degrau {[}ORGANOID{]} pelo protocolo congelado (F1--F10;
GATE-F com assinatura de PI; estimador θ\_obs com scorer cego; plano
estatístico Welch/Holm; kill-switches por braço e programa); (2) degrau
{[}MOUSE{]} sob as regras de pivot pré-declaradas; (3) degrau
{[}HUMAN{]} exclusivamente pela via compassiva E200K com DSMB, endpoints
de contenção/desaceleração, comunicação responsável às famílias e
equidade de acesso à população-âncora; (4) em todos os degraus:
pré-registro antes do dado, publicação do resultado negativo quando
ocorrer, e rótulo de tier em toda saída. Nenhum atalho fora desta
sequência produz resultado aplicável --- é o que a metodologia exige e o
guardião atesta.

\emph{(adendo filosófico REALIZADO: a reflexão etrização/átomo da autora
--- registrada no corpus externo e integrada via Nota à Leitura, §3.1
nome-do-método, Cap. 5.3 declaração mandatória e Cap. 7 síntese --- É o
adendo; commit 83bbf51)}

\hypertarget{capuxedtulo-6-discussuxe3o}{%
\section{CAPÍTULO 6 --- DISCUSSÃO}\label{capuxedtulo-6-discussuxe3o}}

Este capítulo pesa os resultados: o que significam, o que não
significam, e como se posicionam entre a promessa e o limite.

\hypertarget{promessas-e-limites-desta-parte-2}{%
\subsubsection{6.1 Promessas e limites desta Parte
2}\label{promessas-e-limites-desta-parte-2}}

Promete: a etrização como método completo, replicável e em parte já
demonstrado (M1 executado; M4 piloto) --- continuar pesquisa por
simulação hoje, como a física, sem parar à espera de cada confirmação. E
declara: os resultados daqui \textbf{são de simulação e assim permanecem
rotulados}; a aplicação/validação em ambiente real não é prometida nesta
Parte --- se dados reais futuros forem análogos aos simulados,
\textbf{os passos seguintes já estão avançados} (P6)
{[}C054\textbar E009, E010, E031, E032, E033, E007{]}. \textbf{Não
promete}: seleção de parceiro (execução externa), dado {[}ORGANOID{]}+
(não existe ainda --- e é rotulado como tal), nem qualquer claim clínica
(escada de tiers: nenhum degrau empresta a autoridade do seguinte).
Limites declarados: pesos e scores são julgamento estruturado
single-rater; a calibração do estimador é na grade atual (refinamento
futuro = grade mais fina, pré-declarado); PubMed-direto e identificação
R8 pendem como conformidade do piloto.

\hypertarget{discussuxe3o-o-que-os-resultados-significam-e-o-que-nuxe3o-significam}{%
\subsubsection{6.2 Discussão: o que os resultados significam e o que não
significam}\label{discussuxe3o-o-que-os-resultados-significam-e-o-que-nuxe3o-significam}}

Os resultados {[}SIM{]} significam: (i) existe um limiar adimensional
(θ*=0,333) acima do qual o modelo contém a frente --- quantidade
comparável entre subtipos e, no futuro, entre meios
{[}C038\textbar E032{]}; (ii) o design é quantitativo (anel 8--12 mm;
malha ≥5×; redose ≤7 d) --- nenhum candidato anterior o teve
{[}C033\textbar E030{]}; (iii) o método é falseável em três frentes
pré-registradas (θ medido; C50 invariância; dose-resposta do A6
discriminando a forma funcional) {[}C051\textbar E032, E033{]}. Não
significam: contenção em tecido humano medido (isso é {[}ORGANOID{]}+,
inexistente); eficácia clínica; nem substituição do laboratório (§3.3).
A contribuição de maior alcance é a demonstração de que a própria
continuidade da pesquisa pode ser metodologicamente rigorosa sem parar à
espera de confirmação --- a etrização executada de ponta a ponta com
revisão hostil de máquina e validação expressa da autora
{[}C054\textbar E009, E010, E031, E032, E033, E007{]}.

\hypertarget{capuxedtulo-7-conclusuxf5es-por-objetivo}{%
\section{CAPÍTULO 7 --- CONCLUSÕES POR
OBJETIVO}\label{capuxedtulo-7-conclusuxf5es-por-objetivo}}

\textbf{OE1 (nomear e formalizar a etrização) --- alcançado.} etrização
definida com passos P0--P6 e garantias por passo; check BLOCKED próprio
no guardião (R3-BASE-VALIDADE) garante permanência {[}C054\textbar E009,
E010, E031, E032, E033, E007{]}.

\textbf{OE2 (Base de Validade) --- alcançado.} Tríade declarada +
linhagem completa dos seis dados
(quem→espécie→cruzamento→código→parametrização→resultado), com
referências sólidas por linha {[}C046\textbar E032, E033{]}.

\textbf{OE3 (resultados como achados {[}SIM{]}) --- alcançado.} Limiar
θ*=0,333; regras 1--3; sensibilidade com predição discriminadora;
probabilístico de duas lentes; estimador calibrado na fronteira (bias
−0,008) com a v1.1 rejeitada honestamente; tabela-decisão derivada
{[}C038\textbar E032{]} {[}C051\textbar E032, E033{]}
{[}C052\textbar E032, E033{]}.

\textbf{OE4 (continuidade documentada como método, não executada) ---
alcançado.} Loop de re-parametrização anti-hindsight; seleção de
parceiro SLR-análogo (query bank pré-registrada; PubMed-direto conforme
Q1=23/Q2=2; piloto não-decisório); freeze F1--F10 dormente
{[}C053\textbar E033{]}.

\textbf{OE5 (revisão hostil + validação expressa) --- alcançado.} Gates
R0--R3 nos dois perfis, 0 BLOCKED/0 AMEND; AST executável 6/6; validação
expressa da autora registrada em commit (merge do PR \#2).

\textbf{Hipóteses.} H1 corroborada no nível documental (o mapeamento
análogo §3.4 permite leitura por examinador experimental; a banca
decidirá em grau superior). H2 corroborada pelas predições travadas
desde o release v1.0 e pela colheita que as refinou sem movê-las
{[}C051\textbar E032, E033{]}. H3 corroborada pelo posicionamento
estrutural vs in-silico trials (Cap. 2 §2.3).

\textbf{Síntese final.} Esta tese é uma \textbf{etrização} ---
realizada, rotulada e auditável: o método da etrização produziu-a com
garantias (P0--P6), a partir exclusivamente de dados reais publicados, e
a entrega como prognóstico falseável pronto para ser confrontado quando
(e se) o dado real chegar. A tese está realizada nos termos declarados:
pesquisa continuada por simulação parametrizada sobre dados reais, com
prognósticos travados antes de qualquer medição e continuidade futura
documentada como método. Se dados reais vierem a ser análogos aos
simulados, os passos seguintes já estão avançados (P6) --- antecipação
bancada, nunca pendente {[}C054\textbar E009, E010, E031, E032, E033,
E007{]}.

\textbf{Fecho epistemológico.} Esta tese não apenas demonstra a
etrização --- ela \textbf{legitima uma nova forma de fazer ciência na
saúde}. A questão filosófica que a motivou é simples: \emph{o que uma
pesquisadora pode legítimamente produzir quando a doença mata em meses,
o laboratório custa fortunas, e os dados publicados por milhares de
pesquisadores estão disponíveis?} A resposta desta tese:
\textbf{prognóstico quantitativo, falseável e disciplinado por
critérios} --- com nome (etrização), método (P0-P6), garantias (5
requisitos) e honestidade sobre o que é e o que não é. Como o átomo
tornou pesquisável o invisível, a etrização torna pesquisável o
ainda-não-medido.

\hypertarget{apuxeandice-a-inventuxe1rio-dos-artefatos-da-parte-2-verificuxe1vel-em-disco}{%
\section{APÊNDICE A --- Inventário dos artefatos da Parte 2 (verificável
em
disco)}\label{apuxeandice-a-inventuxe1rio-dos-artefatos-da-parte-2-verificuxe1vel-em-disco}}

\begin{longtable}[]{@{}
  >{\raggedright\arraybackslash}p{(\columnwidth - 4\tabcolsep) * \real{0.3333}}
  >{\raggedright\arraybackslash}p{(\columnwidth - 4\tabcolsep) * \real{0.3333}}
  >{\raggedright\arraybackslash}p{(\columnwidth - 4\tabcolsep) * \real{0.3333}}@{}}
\toprule\noalign{}
\begin{minipage}[b]{\linewidth}\raggedright
Artefato
\end{minipage} & \begin{minipage}[b]{\linewidth}\raggedright
O que é
\end{minipage} & \begin{minipage}[b]{\linewidth}\raggedright
Estado
\end{minipage} \\
\midrule\noalign{}
\endhead
\bottomrule\noalign{}
\endlastfoot
\texttt{experiments/part2\_theta\_obs\_v1.py} +
\texttt{part2\_results/part2\_theta\_obs\_v1.json} & estimador v1-NN:
grade κ 1,5--8 + calibração unitária (veredito ADEQUADO pré-declarado) &
executado \\
\texttt{experiments/part2\_theta\_obs\_pooled.py} + \texttt{.json} &
regime declarado §2.7: mediana por braço n=8 (bias −0,008 na fronteira;
modal 69\%) & executado \\
\texttt{experiments/part2\_theta\_obs\_v11.py} + \texttt{.json} &
estimador interpolado IDW-2 --- \textbf{testado e REJEITADO} (bias
anti-conservador na fronteira; cobertura quebra em κ=8); JSON regenerado
por script após truncamento da primeira execução (nota de integridade
in-file) & executado (rejeitado) \\
\texttt{experiments/part2\_results/part2\_derived\_summary.json} &
tabela-decisão derivada (κ↔θ↔R↔biomassa↔margem) --- zero novas
simulações & derivado \\
\texttt{experiments/part2\_results/partner\_selection\_log.md} v0.2 &
aplicação-piloto do método M4 (9 registros → 8 grupos → 2+1+3+2) ---
não-decisório & piloto \\
\texttt{experiments/G0\_EXECUTION\_FREEZE\_CHECKLIST.md} & F1--F10 +
GATE-F (M2 --- dormant por reenquadramento) & especificado \\
\texttt{experiments/REPARAM\_LOOP.md} & regra anti-hindsight de
realimentação (M3) & especificado \\
\texttt{experiments/PARTNER\_SELECTION\_PROTOCOL.md} v2.1 & método M4
(query bank + I/X + pesos + fluxo) & método \\
\texttt{paper/guardian/guardian.py} (\texttt{-\/-profile\ part2}) & gate
da Parte 2 (inclui R3-BASE-VALIDADE BLOCKED) & operante \\
\end{longtable}

\emph{Nota de integridade da auditoria 27/08 (tarde): dois defeitos
encontrados e corrigidos no próprio ato --- JSON da v11 truncado por
erro de serialização (regenerado por script salvo, seed idêntica,
veredito idêntico) e sinal do bias pooled invertido na prosa (−0,008
conforme o JSON; corrigido nos três documentos). Ambos registrados aqui
porque a auditoria que não publica suas correções não é auditoria.}

\hypertarget{referuxeancias-e-concorduxe2ncia}{%
\subsection{REFERÊNCIAS E
CONCORDÂNCIA}\label{referuxeancias-e-concorduxe2ncia}}

\hypertarget{referuxeancias-38-fontes-do-registro-geradas-do-manifest}{%
\subsubsection{Referências (38 fontes do registro, geradas do
manifest)}\label{referuxeancias-38-fontes-do-registro-geradas-do-manifest}}

{[}1{]} ASANTE, E.A. et al.~A naturally occurring variant of the human
prion protein completely prevents prion disease. 2015.
doi:10.1038/nature14510. {[}2{]} MEAD, S. et al.~A novel protective
prion protein variant that colocalizes with kuru exposure. 2009.
doi:10.1056/NEJMoa0809716. {[}3{]} GATDULA, J.R. et al.~Leveraging the
dominant-negative effect of the kuru-protective G127V prion protein
variant as a novel therapeu. 2026. PMID 41757113. {[}4{]} ZERBES, T. et
al.~A self-complementary recombinant adeno-associated virus vector
coding for an anchorless prion protein carrying. 2026. . {[}5{]} HOSSZU,
L.P. et al.~Structural effects of the highly protective V127
polymorphism on human prion protein. 2020.
doi:10.1038/s42003-020-01126-6. {[}6{]} ZHENG, Z. et al.~Structural
basis for the complete resistance of the human prion protein mutant
G127V to prion disease. 2018. doi:10.1038/s41598-018-31394-6. {[}7{]}
GROVEMAN, B.R. et al.~Sporadic Creutzfeldt-Jakob disease prion infection
of human cerebral organoids. 2019. doi:10.1186/s40478-019-0742-2.
{[}8{]} GROVEMAN, B.R. et al.~Human cerebral organoids as a therapeutic
drug screening model for Creutzfeldt-Jakob disease. 2021.
doi:10.1038/s41598-021-84689-6. {[}9{]} FORNARA, B. et al.~The dynamics
of prion spreading is governed by the interplay between the
non-linearities of tissue response an. 2024. PMID 39717079. {[}10{]}
THORNE, R.G. et al.~In vivo diffusion analysis with quantum dots and
dextrans predicts extracellular space and tortuosity in brain. 2006.
doi:10.1073/pnas.0509425103. {[}11{]} MASEL, J. et al.~Quantifying the
kinetic parameters of prion replication. 1999.
doi:10.1016/S0301-4622(99)00004-3. {[}12{]} WILLIAMS, K. et al.~Neural
cell engraftment therapy for sporadic Creutzfeldt-Jakob disease restores
neuroelectrophysiological para. 2023. doi:10.1186/s13287-023-03591-2.
{[}13{]} RELANO-GINES, A. et al.~Prion replication occurs in endogenous
adult neural stem cells and alters their neuronal fate. 2013.
doi:10.1371/journal.ppat.1003485. {[}14{]} GINHOUX, F. et al.~Fate
mapping analysis reveals that hematopoietic cells of yolk-sacil origin
give rise to microglia. 2010. doi:10.1126/science.1194637. {[}15{]}
SORRELLS, S.F. et al.~Human hippocampal neurogenesis drops sharply in
children to undetectable levels in adults. 2018.
doi:10.1038/s41586-018-0336-4. {[}16{]} ABUD, E.M. et al.~iPSC-derived
human microglia-like cells to study neurological diseases. 2017. PMID
28426964. {[}17{]} HAN, X. et al.~Generation of hypoimmunogenic human
pluripotent stem cells. 2019. doi:10.1073/pnas.1902566116. {[}18{]} HU,
X. et al.~Hypoimmune induced pluripotent stem cells survive long-term in
fully immunocompetent allogeneic rhesus macaque. 2024.
doi:10.1038/s41587-023-01784-x. {[}19{]} XUE, Y. et al.~Lipid
nanoparticles enhance mRNA delivery to the central nervous system upon
intrathecal injection. 2025. PMID 40317512. {[}20{]} LIANG, Y. et
al.~The survival of engrafted neural stem cells within hyaluronic acid
hydrogels. 2013. PMID 23623429. {[}21{]} SHAH, S.Z. et al.~Early
minocycline and late FK506 treatment improves survival\ldots{} in
prion-infected hamsters (RETRACTED). 2017.
doi:10.1007/s13311-020-00909-3. {[}22{]} CHENG, S. et al.~Minocycline
reduces neuroinflammation but does not improve survival in
prion-infected mice. 2015. doi:10.1038/srep10535. {[}23{]} GENTILE, J.E.
et al.~Evidence that minocycline treatment confounds neurofilament light
chain biomarker interpretation. 2024.
https://www.ukdri.ac.uk/publications/evidence-minocycline-treatment-confounds-interpretation-neurofilament-biomarker.
{[}24{]} SMID, J. et al.~Creutzfeldt-Jakob disease associated with a
missense mutation at codon 200 of the prion protein gene in Brazil.
2007.
https://www.demneuropsy.com.br/article/creutzfeldt-jakob-disease-associated-with-a-missense-mutation-at-codon-200-of-the-prion-protein-gene-in-brazil/.
{[}25{]} STOPSCHINSKI, B.E. et al.~Prion-like mechanisms in
neurodegenerative disease. 2017. doi:10.1016/S1474-4422(17)30370-6.
{[}26{]} JUCKER, M. et al.~Propagation and spread of pathogenic protein
assemblies in neurodegenerative diseases. 2018.
doi:10.1038/s41586-018-0344-4. {[}27{]} FDA, (.F. et al.~FDA grants
accelerated approval of tofersen for SOD1-ALS (press release/decision
summary). 2023.
https://www.fda.gov/news-events/press-announcements/fda-grants-accelerated-approval-first-treatment-als-patients-rare-genetic-form-disease.
{[}28{]} FDA, (.F. et al.~FDA approval of nusinersen (Spinraza) ---
regulatory record. 2016.
https://www.fda.gov/vaccines-blood-biologics/approved-blood-products/spinraza-nusinersen.
{[}29{]} BENGTSSON, S. et al.~Clinical trial of stem-cell derived
dopaminergic progenitor transplantation in Parkinson's disease
(feasibilit. 2026. https://www.newscientist.com/article/\ldots. {[}30{]}
OPEN, P.\&. et al.~WS-7: ADR transport solver --- self-tested design
rules (this program). 2026. Disponível em:
https://github.com/camillanapoles/quest003-prion-v127. {[}31{]} OPEN,
P.\&. et al.~WS-8: hierarchical Bayesian calibration over structural
analogues (this program). 2026. Disponível em:
https://github.com/camillanapoles/quest003-prion-v127. {[}32{]} OPEN,
P.\&. et al.~WS-9: humanized in-silico infection model with V127 capping
(this program). 2026. Disponível em:
https://github.com/camillanapoles/quest003-prion-v127. {[}33{]} OPEN,
P.\&. et al.~Quest 003 repository --- timestamped pre-registrations and
audit trail (this program). 2026. Disponível em:
https://github.com/camillanapoles/quest003-prion-v127. {[}34{]}
GESCHWIND, M.D.~et al.~Quinacrine treatment trial for sporadic
Creutzfeldt-Jakob disease (Class I: no survival benefit). 2013. PMID
24122181. {[}35{]} HAIK, S. et al.~Doxycycline in Creutzfeldt-Jakob
disease: a phase 2, randomised, double-blind, placebo-controlled trial
(no si. 2014. PMID 24411709. {[}36{]} NEWMAN, P.K. et al.~Postmortem
findings in a case of variant CJD treated with intraventricular pentosan
polysulfate (iPPS): biolog. 2014. PMID 24554103. {[}37{]} OTTO, M. et
al.~Efficacy of flupirtine on cognitive function in patients with CJD: a
double-blind placebo-controlled study (co. 2004.
doi:10.1212/01.WNL.0000113764.35026.ef. {[}38{]} MEAD, S. et al.~Prion
protein monoclonal antibody (PRN100) therapy for Creutzfeldt-Jakob
disease: evaluation of a first-in-hum. 2022. PMID 35305340.

\hypertarget{fontes-complementares-skill-scout-verificauxe7uxe3o-completa-pendente-snippet-level}{%
\subsubsection{Fontes complementares (skill-scout; verificação completa
pendente ---
snippet-level)}\label{fontes-complementares-skill-scout-verificauxe7uxe3o-completa-pendente-snippet-level}}

\begin{itemize}
\tightlist
\item
  CORTÉS-RÍOS, J. et al.~A step-by-step workflow for performing in
  silico clinical trials. 2025. PMCID PMC12706418. \emph{(abertura
  completa pendente antes de registro E)}
\item
  An open source statistical web application for validation of in-silico
  trials and virtual cohorts. 2025. doi:10.1038/s41598-025-99720-3.
  \emph{(idem)}
\end{itemize}

\hypertarget{concorduxe2ncia-claims-referuxeancias-ruxe9gua-da-autora-claim-sem-referuxeancia-uxe9-inaceituxe1vel}{%
\subsubsection{Concordância claims → referências (régua da autora: claim
sem referência é
inaceitável)}\label{concorduxe2ncia-claims-referuxeancias-ruxe9gua-da-autora-claim-sem-referuxeancia-uxe9-inaceituxe1vel}}

\begin{longtable}[]{@{}
  >{\raggedright\arraybackslash}p{(\columnwidth - 4\tabcolsep) * \real{0.3333}}
  >{\raggedright\arraybackslash}p{(\columnwidth - 4\tabcolsep) * \real{0.3333}}
  >{\raggedright\arraybackslash}p{(\columnwidth - 4\tabcolsep) * \real{0.3333}}@{}}
\toprule\noalign{}
\begin{minipage}[b]{\linewidth}\raggedright
Claim
\end{minipage} & \begin{minipage}[b]{\linewidth}\raggedright
Evidências
\end{minipage} & \begin{minipage}[b]{\linewidth}\raggedright
Referências
\end{minipage} \\
\midrule\noalign{}
\endhead
\bottomrule\noalign{}
\endlastfoot
C001 & E001 & {[}1{]} \\
C002 & E001 & {[}1{]} \\
C003 & E001, E003, E005 & {[}1{]}, {[}3{]}, {[}5{]} \\
C004 & E002 & {[}2{]} \\
C005 & E003 & {[}3{]} \\
C006 & E003 & {[}3{]} \\
C007 & E004 & {[}4{]} \\
C008 & E005, E006 & {[}5{]}, {[}6{]} \\
C009 & E007 & {[}7{]} \\
C010 & E007 & {[}7{]} \\
C011 & E007 & {[}7{]} \\
C012 & E008 & {[}8{]} \\
C013 & E009 & {[}9{]} \\
C014 & E010, E030 & {[}10{]}, {[}30{]} \\
C015 & E011, E030 & {[}11{]}, {[}30{]} \\
C016 & E012 & {[}12{]} \\
C017 & E013 & {[}13{]} \\
C018 & E014 & {[}14{]} \\
C019 & E015 & {[}15{]} \\
C020 & E016 & {[}16{]} \\
C021 & E017, E018 & {[}17{]}, {[}18{]} \\
C022 & E019 & {[}19{]} \\
C023 & E020 & {[}20{]} \\
C024 & E021 & {[}21{]} \\
C025 & E022 & {[}22{]} \\
C026 & E023 & {[}23{]} \\
C027 & E024 & {[}24{]} \\
C028 & E025, E026 & {[}25{]}, {[}26{]} \\
C029 & E027 & {[}27{]} \\
C030 & E028 & {[}28{]} \\
C031 & E029 & {[}29{]} \\
C032 & E030 & {[}30{]} \\
C033 & E030 & {[}30{]} \\
C034 & E030, E020 & {[}30{]}, {[}20{]} \\
C035 & E030, E019 & {[}30{]}, {[}19{]} \\
C036 & E031 & {[}31{]} \\
C037 & E032, E007 & {[}32{]}, {[}7{]} \\
C038 & E032 & {[}32{]} \\
C039 & E032, E007 & {[}32{]}, {[}7{]} \\
C040 & E033 & {[}33{]} \\
C041 & E033 & {[}33{]} \\
C042 & E030 & {[}30{]} \\
C043 & E032, E009 & {[}32{]}, {[}9{]} \\
C044 & E032 & {[}32{]} \\
C045 & E033 & {[}33{]} \\
C046 & E032, E009, E007, E033 & {[}32{]}, {[}9{]}, {[}7{]}, {[}33{]} \\
C047 & E032, E033 & {[}32{]}, {[}33{]} \\
C048 & E009, E010, E030, E031, E032, E033 & {[}9{]}, {[}10{]}, {[}30{]},
{[}31{]}, {[}32{]}, {[}33{]} \\
C049 & E032, E033 & {[}32{]}, {[}33{]} \\
C050 & E021, E022, E034, E035, E036, E037, E038 & {[}21{]}, {[}22{]},
{[}34{]}, {[}35{]}, {[}36{]}, {[}37{]}, {[}38{]} \\
C051 & E032, E033 & {[}32{]}, {[}33{]} \\
C052 & E032, E033 & {[}32{]}, {[}33{]} \\
C053 & E033 & {[}33{]} \\
C054 & E009, E010, E031, E032, E033, E007 & {[}9{]}, {[}10{]}, {[}31{]},
{[}32{]}, {[}33{]}, {[}7{]} \\
\end{longtable}

\hypertarget{parte-2-sem-ano-o-resultado-define-as-fases-suxe3o-estimativas.-toda-cifra-deste-manuscrito-vem-de-json-arquivado-ou-do-registro-e-regra-nunca-digitar-valor.-gated-guardian-r0r3-0-blocked-exigido.}{%
\subsection{\texorpdfstring{\emph{Parte 2 \{SEM ANO\} --- o resultado
define; as fases são estimativas. Toda cifra deste manuscrito vem de
JSON arquivado ou do registro E (regra: nunca digitar valor). Gated:
guardian R0--R3, 0 BLOCKED
exigido.}}{Parte 2 \{SEM ANO\} --- o resultado define; as fases são estimativas. Toda cifra deste manuscrito vem de JSON arquivado ou do registro E (regra: nunca digitar valor). Gated: guardian R0--R3, 0 BLOCKED exigido.}}\label{parte-2-sem-ano-o-resultado-define-as-fases-suxe3o-estimativas.-toda-cifra-deste-manuscrito-vem-de-json-arquivado-ou-do-registro-e-regra-nunca-digitar-valor.-gated-guardian-r0r3-0-blocked-exigido.}}

\emph{Parte 2 \{SEM ANO\} --- o resultado define; as fases são
estimativas. Toda cifra vem de JSON arquivado ou do registro E (regra:
nunca digitar valor). Gated: guardian R0--R3 perfil part2, 0 BLOCKED
exigido. PT é o mestre; EN companion: manuscript\_Parte2\_v1\_EN.md.}

\hypertarget{apuxeandice-b-mapa-da-luxf3gica-completa-de-onde-veio-cada-informauxe7uxe3o-e-como-se-conecta}{%
\section{APÊNDICE B --- MAPA DA LÓGICA COMPLETA (de onde veio cada
informação e como se
conecta)}\label{apuxeandice-b-mapa-da-luxf3gica-completa-de-onde-veio-cada-informauxe7uxe3o-e-como-se-conecta}}

\hypertarget{b.1-fluxo-de-informauxe7uxe3o-da-literatura-uxe0-conclusuxe3o}{%
\subsection{B.1 Fluxo de informação: da literatura à
conclusão}\label{b.1-fluxo-de-informauxe7uxe3o-da-literatura-uxe0-conclusuxe3o}}

\begin{verbatim}
CAMADA 0: FONTE PRIMÁRIA (papers publicados)
═══════════════════════════════════════════════════
  Asante 2015 [E001]     Mead 2009 [E002]     Groveman 2019 [E007]
  Nature, peer-reviewed   NEJM, peer-reviewed   Acta Neuropathol, peer-reviewed
  "V127 resiste a príons" "G127V seleção kuru"  "Organoides infectam com sCJD"
         │                       │                       │
         ▼                       ▼                       ▼
CAMADA 1: VERIFICAÇÃO (auditoria de fonte)
═══════════════════════════════════════════════════
  Cada fonte foi ABERTA (URL/DOI confirmado)
  Identificadores verificados → E-registry
  Status: verified (quem abriu, quando, como)
         │
         ▼
CAMADA 2: CLAIM REGISTRY (afirmações extraídas)
═══════════════════════════════════════════════════
  C001: "V127/V127 resiste completamente"      [E001]
  C003: "DN dose-dependente"                   [E001,E003,E005]
  C005: "anchorless trans DN"                  [E003]
  C010: "âncoras de relógio: 25-28/35/169 dpi" [E007]
  C011: "títulos MV2 vs MV1 @169dpi"           [E007]
         │
         ▼
CAMADA 3: PARÂMETROS (números com proveniência)
═══════════════════════════════════════════════════
  α=0.20, λ=1.8     ← Thorne & Nicholson 2006 [E010] (in-vivo, humano)
  D₀=1.25e-10 m²/s  ← Stokes-Einstein (R_h≈2.5nm, ~30kDa)
  D_eff=3.86e-11    ← D₀/λ² (derivado)
  k_eff=1e-6–1e-5   ← Masel 1999 [E011] (polimerização nucleada)
  t_dupl=12.1 d     ← 134d/log₂(2130) [E007] (derivado de Groveman)
  1 unidade=144 d   ← t_dupl/t_dupl_sim [E032] (derivado)
         │
         ▼
CAMADA 4: MODELO MATEMÁTICO (equações com parâmetros)
═══════════════════════════════════════════════════
  ∂(αc)/∂t = ∇·(D_eff∇c) − ∇·(vc) + S(x) − k_eff·c
  → solver WS-7 (volumes finitos, 192², auto-testado)

  freeS(r) = (1+κ·c_V127(r))⁻²
  → modelo WS-9 (campo-médio, 96², humanizado)

  θ ≡ (1+κ·c_pico)⁻¹
  → definição formal (a fração de replicação que resta)
         │
         ▼
CAMADA 5: EXECUÇÃO (runs arquivados e reproduzidos)
═══════════════════════════════════════════════════
  WS-7 self-tests: massa=100%, Thiele=0.5% ✓
  WS-9 T1/T2/T3: todos PASS ✓
  Reprodução: executado em 2 ambientes independentes
  → resultados idênticos (hash confirmado)
         │
         ▼
CAMADA 6: RESULTADOS [SIM]
═══════════════════════════════════════════════════
  θ*=0.333          ← do sweep κ (limiar de contenção)
  Regras 1-3        ← do solver WS-7 (anel/malha/redose)
  Discriminadora    ← do sweep expoente (C50 insensível)
  Estimador θ_obs   ← da calibração sim-a-sim
         │
         ▼
CAMADA 7: PROGNÓSTICO TRAVADO
═══════════════════════════════════════════════════
  Release v1.0 (timestamp): "θ<0.33 ⇒ contenção"
  → ANTES de qualquer dado experimental
  → ANTI-HINDSIGHT por construção
         │
         ▼
CAMADA 8: CONCLUSÕES DA TESE
═══════════════════════════════════════════════════
  A tese entrega: método (etrização P0-P6) + prognóstico (θ*, regras)
  + dupla potencialidade (cura + restauração [C016])
  + ressalva (não aplicar em humanos sem validação experimental)
\end{verbatim}

\hypertarget{b.2-decisuxf5es-chave-por-que-cada-escolha-foi-feita}{%
\subsection{B.2 Decisões-chave: por que cada escolha foi
feita}\label{b.2-decisuxf5es-chave-por-que-cada-escolha-foi-feita}}

\begin{longtable}[]{@{}
  >{\raggedright\arraybackslash}p{(\columnwidth - 4\tabcolsep) * \real{0.3333}}
  >{\raggedright\arraybackslash}p{(\columnwidth - 4\tabcolsep) * \real{0.3333}}
  >{\raggedright\arraybackslash}p{(\columnwidth - 4\tabcolsep) * \real{0.3333}}@{}}
\toprule\noalign{}
\begin{minipage}[b]{\linewidth}\raggedright
Decisão
\end{minipage} & \begin{minipage}[b]{\linewidth}\raggedright
Alternativa considerada
\end{minipage} & \begin{minipage}[b]{\linewidth}\raggedright
Por que esta foi escolhida
\end{minipage} \\
\midrule\noalign{}
\endhead
\bottomrule\noalign{}
\endlastfoot
V127 como agente & ASO (reduzir PrP) & DN preserva PrP nativo (função
sináptica); ASO depleta \\
Anchorless (sem GPI) & V127 de membrana & Gatdula 2026: funciona em
trans; contém vizinhas; é difusível \\
Modelo determinístico & Manter Gillespie estocástico & Mean-field
permite sweep sistemático; estocástico é barulhento \\
Reescala global de tempo & Re-ajustar cada taxa individualmente & Só um
grau de liberdade calibrável (relógio de Groveman); re-ajustar 6+ taxas
seria sobreajuste \\
freeS quadrático & freeS linear & Assumimos conversão de 2
participantes; a alternativa é TESTÁVEL (predição discriminadora) \\
θ como limiar & R (raio) como limiar & θ é adimensional (comparável
entre subtipos); R satura \\
P(G0)=36.6\% & P(G0)=62\% estrutural & Duas lentes mantidas separadas: a
empírica (36.6\%) é conservadora; a estrutural (62\%) é arquivada sem
mescla \\
\end{longtable}

\hypertarget{b.3-rejeiuxe7uxf5es-documentadas-o-que-nuxe3o-funcionou-e-por-quuxea}{%
\subsection{B.3 Rejeições documentadas (o que NÃO funcionou e por
quê)}\label{b.3-rejeiuxe7uxf5es-documentadas-o-que-nuxe3o-funcionou-e-por-quuxea}}

\begin{longtable}[]{@{}
  >{\raggedright\arraybackslash}p{(\columnwidth - 4\tabcolsep) * \real{0.3333}}
  >{\raggedright\arraybackslash}p{(\columnwidth - 4\tabcolsep) * \real{0.3333}}
  >{\raggedright\arraybackslash}p{(\columnwidth - 4\tabcolsep) * \real{0.3333}}@{}}
\toprule\noalign{}
\begin{minipage}[b]{\linewidth}\raggedright
O que foi rejeitado
\end{minipage} & \begin{minipage}[b]{\linewidth}\raggedright
Por quê
\end{minipage} & \begin{minipage}[b]{\linewidth}\raggedright
Onde está
\end{minipage} \\
\midrule\noalign{}
\endhead
\bottomrule\noalign{}
\endlastfoot
v1.1-IDW (estimador interpolado) & PIOROU a fronteira de decisão (bias
−0.037, direção anti-conservadora) & part2\_theta\_obs\_v11.json \\
SVZ como fábrica central & Nicho humano quiescente + replica príon
(refutação dupla) & R-1 no canon \\
Bialélico como suficiente & Heterozigoto é infectável por vCJD & R-2 no
canon \\
Ensaio retraído (minociclina/FK506) & Retraído em 2020; ainda circulava
como evidência & Excluído por regra \\
NSC→micróglia via transdiferenciação & Ginhoux 2010: linhagem impossível
& Corrigido → co-enxerto iMG \\
\end{longtable}

\hypertarget{b.4-glossuxe1rio-contextual-todo-termo-tuxe9cnico-definido-onde-aparece}{%
\subsection{B.4 Glossário contextual (todo termo técnico definido onde
aparece)}\label{b.4-glossuxe1rio-contextual-todo-termo-tuxe9cnico-definido-onde-aparece}}

\begin{longtable}[]{@{}
  >{\raggedright\arraybackslash}p{(\columnwidth - 4\tabcolsep) * \real{0.3333}}
  >{\raggedright\arraybackslash}p{(\columnwidth - 4\tabcolsep) * \real{0.3333}}
  >{\raggedright\arraybackslash}p{(\columnwidth - 4\tabcolsep) * \real{0.3333}}@{}}
\toprule\noalign{}
\begin{minipage}[b]{\linewidth}\raggedright
Termo
\end{minipage} & \begin{minipage}[b]{\linewidth}\raggedright
Definição na tese
\end{minipage} & \begin{minipage}[b]{\linewidth}\raggedright
Primeira aparição
\end{minipage} \\
\midrule\noalign{}
\endhead
\bottomrule\noalign{}
\endlastfoot
etrização & método de continuar pesquisa por simulação com dados reais
{[}C054{]} & Nota à Leitura \\
θ & fração de replicação remanescente no pico secretor (1+κc)⁻¹ &
§3.2 \\
κ & força de capping (capacidade do V127ΔGPI de bloquear conversão) &
§3.2 \\
G0-sim & gate computacional executado e aprovado (T1/T2/T3) & §1.7 \\
G0-wet & gate de organoide especificado mas não executado (dormant) &
§4.3 \\
tier {[}SIM{]} & dado produzido por simulação parametrizada & toda
saída \\
tier {[}ORGANOID{]} & dado de laboratório com organoide (inexistente
ainda) & §5.3 \\
DN & dominante-negativo: V127 bloqueia conversão de PrP normal &
Cap.2 \\
θ* & limiar de contenção: acima dele a frente é contida (0.333) &
§4.1 \\
freeS & fração de substrato livre após capping: (1+κc)⁻² & §3.2 \\
\end{longtable}

\hypertarget{b.5-prejulgando-objeuxe7uxf5es-da-banca}{%
\subsection{B.5 Prejulgando objeções da
banca}\label{b.5-prejulgando-objeuxe7uxf5es-da-banca}}

\begin{longtable}[]{@{}
  >{\raggedright\arraybackslash}p{(\columnwidth - 2\tabcolsep) * \real{0.5000}}
  >{\raggedright\arraybackslash}p{(\columnwidth - 2\tabcolsep) * \real{0.5000}}@{}}
\toprule\noalign{}
\begin{minipage}[b]{\linewidth}\raggedright
Objeção provável
\end{minipage} & \begin{minipage}[b]{\linewidth}\raggedright
Resposta que está na tese
\end{minipage} \\
\midrule\noalign{}
\endhead
\bottomrule\noalign{}
\endlastfoot
``Você não tem dados experimentais'' & Esta é uma tese de METODOLOGIA
com achados de simulação auditada; a descoberta biológica virá do
G0-wet; a tese entrega método + prognóstico travado + dupla
potencialidade declarada com ressalva \\
``A simulação não substitui o laboratório'' & Correto --- e a tese
DECLARA isto na Base de Validade (§3.3): ``não substitui; estimula P\&D
ágil'' \\
``Como sei que os números são reais?'' & Cada número tem fonte E-ID com
hash; cada fonte foi aberta e verificada; a reprodução em 2 ambientes
confirmou os valores; a tabela de concordância claims→refs (no fim)
permite auditoria linha por linha \\
``Por que confiar no θ*=0.333?'' & É uma PREDIÇÃO, não uma descoberta;
foi travada antes de qualquer dado; o estimador θ\_obs está calibrado
para testá-la; o modelo tem 3 frentes de falseabilidade (θ, C50,
dose-resposta) \\
``Isto é filosofia, não ciência'' & É metodologia com 38 fontes
verificadas, 54 claims com hash, 48 fatos numéricos, equações com
proveniência, simulação reproduzida, revisão hostil de máquina, e
pré-registro timestampado \\
``Como se aplica a Alzheimer?'' & §5.2-bis: exemplo concreto P0→P5; MAS:
para Aβ não existe variante protetora natural (V127 era do kuru); o
agente teria de ser engenheirado --- incerteza adicional declarada \\
\end{longtable}

\end{document}
